\begin{quote}
\textbf{Nota sobre esta versión.} Esta es una traducción técnica al español del
paper en inglés \texttt{paper\_v2\_en.md}, mantenida en paralelo al original. La
fuente de verdad para arXiv y para todas las tablas / claims experimentales
sigue siendo la versión en inglés. Se preservan en inglés los términos técnicos
reservados de la literatura (LLM, RAG, HNSW, snapshot cache, key-value, sorted
set, HyperLogLog, NEON SDOT, SIMD, in-process, write-through, backend, daemon,
thread, lock, atomicity, prefill, decode, embedding, retrieval, fork, harness,
benchmark, footprint, ingest, wrapper, etc.) tal como aparecen en el dominio.
\end{quote}

\section{Introducción}\label{introducciuxf3n}

Los modelos de lenguaje cuantizados de menos de 3 GB hoy corren a tasas
interactivas en hardware ARM64 de clase smartphone: Gemma 2 {[}@gemma2{]} y
Gemma 4 {[}@gemma4{]}, Phi-3 mini {[}@phi3{]}, la serie Llama 3 {[}@llama3{]}
(incluyendo Llama 3.2 1B y 3B), MobileLLM {[}@mobilellm{]}, y la serie Qwen 2.5
{[}@qwen25{]} que usamos en §5.9. El stack de runtime --- \texttt{llama.cpp}
{[}@llamacpp{]} con su formato de pesos cuantizados GGUF {[}@gguf{]}, LiteRT-LM
{[}@litertlm{]}, MLC-LLM {[}@mlcllm{]}, ExecuTorch {[}@executorch{]} --- está
maduro y alcanza latencia sub-segundo por token decodificado en handsets de
\$150.

La \textbf{capa de estado} que un agente necesita no ha madurado al mismo ritmo.
Por capa de estado nos referimos al lugar donde el agente guarda lo que ha visto
entre llamadas de inferencia, coordina entre turnos, y sobrevive reinicios de
proceso.

Considere un agente de monitoreo IoT móvil: recolecta lecturas de temperatura,
humedad, presión y batería a intervalos; un LLM procesa lotes en checkpoints y
produce reportes de síntesis. Entre checkpoints el agente necesita (1) un
bounded buffer ordenado en el tiempo sobre las lecturas más recientes, (2)
agregados en running (min/max/avg/count) sin re-escanear historia, (3) conteo
probabilístico para señales de anomalía y cardinalidad, y (4) TTL por campo para
mantener el bloque de contexto dentro del presupuesto del prompt. Cada
requerimiento mapea limpiamente a una estructura de datos específica: bounded
streams, sorted sets y hashes, HyperLogLog, TTL sobre hash fields. Ninguna es
exótica. Ninguna es problema de investigación.

Ningún backend en la plataforma las ofrece como primitivas nativas. En Android e
iOS las opciones embebidas dominantes --- SQLite {[}@sqlite{]}, RocksDB
{[}@rocksdb{]} (un descendiente LSM-tree de LevelDB {[}@leveldb{]}), LMDB
{[}@lmdb{]}, ObjectBox {[}@objectbox4{]}, y el peer embebido analítico DuckDB
{[}@duckdb{]} --- exponen entre dos y cuatro estructuras de datos en total:
filas, key-value pairs, blob stores, ocasionalmente un índice vectorial. Cada
construcción de orden superior (bounded streams, agregados materializados,
conteos aproximados, semántica TTL, similarity search) tiene que ser
re-implementada en código de aplicación sobre esas primitivas, por cada equipo
que las necesita. El código resultante es frágil bajo acceso concurrente. Las
garantías de atomicidad que vienen gratis con un servidor in-process maduro se
pierden silenciosamente. La misma media docena de bugs se redescubre app tras
app.

La aritmética de latencia libera al argumento de la habitual carrera de
microbenchmarks. Incluso el backend embebido más lento en nuestra evaluación ---
SQLite a aproximadamente 3 ms de retrieval en N = 20 000 bajo carga concurrente
en un handset Android de \$150 --- está tres órdenes de magnitud por debajo de
un solo turno de inferencia LLM on-device (1--3 segundos para un modelo de 1--3
miles de millones de parámetros). El retrieval en microsegundos no es lo que
mueve la aguja del producto. Lo que la mueve es qué primitivas expone el
backend, qué tan ergonómicas son de invocar, y cuánto código pegamento escribe
el autor del agente antes de que la capa de estado se comporte.

El eje correcto de comparación para backends de agentes on-device, argumentamos,
es \textbf{superficie primitiva y ergonomía del desarrollador}, no throughput
crudo de microbenchmark. Cuatro piezas de evidencia respaldan esa posición: (a)
el envelope de latencia de Dazzle se ubica cómodamente dentro del piso de ruido
LLM en cada configuración que probamos; (b) primitiva-por-primitiva Dazzle
expone \textbf{diez primitivas nativas contra \(\leq 4\) en cualquier competidor
embebido en las mismas plataformas}; (c) una comparación en líneas de código
para la misma capa de estado del agente; y (d) una demostración end-to-end de
que las primitivas que Dazzle expone habilitan patrones cualitativamente nuevos
--- small-model + on-device retrieval-augmented generation {[}@lewisrag;
@karpukhindpr{]} --- que no son alcanzables cuando la única primitiva disponible
es un store key-value plano o de filas.

\textbf{Contribuciones.}

\begin{enumerate}
\def\labelenumi{\arabic{enumi}.}
\tightlist
\item
  \textbf{Primer port in-process de Valkey a procesos de aplicación Android e
  iOS}: un transporte self-pipe, un snapshot cache con write-through, y un SPSC
  ring buffer en Android reemplazan el TCP loopback de Valkey preservando la
  superficie completa de primitivas y el scripting Lua.
\item
  \textbf{Una contribución de ingeniería medida más dos capas de
  future-proofing} (§3, §5.4). El snapshot cache write-through es la única capa
  que mueve throughput sobre la carga Write-Heavy / Read-Light medida --- la
  ablación factorial 2³ de §5.4 muestra que el worker pool paralelo y el path de
  dispatch por hash-bucket son no-ops sobre esta carga (dentro del ruido
  run-to-run) porque el snapshot cache ya absorbe la contención que ellos fueron
  diseñados para remover. Mantenemos esas dos capas en la arquitectura como
  headroom de keyspace-scaling: bajo benchmarks no corridos en este paper
  (keyspaces mucho más grandes, múltiples threads escritores, fan-in desde
  varias producer queues) esperamos que ganen su costo de vuelta, y los números
  de §5.4 describen las condiciones bajo las que no lo hacen. El auto-mirror
  post-EVAL que mantiene visible al cache el estado mutado por Lua, y el
  transporte in-process en sí, son las dos contribuciones de ingeniería
  adicionales de interés independiente documentadas en §3.
\item
  \textbf{Caracterización del envelope de latencia} sobre doce backends, con N
  de 200 a 20 000, en un handset Android de presupuesto (Moto G35 5G, Unisoc
  T760) y un iPhone 12 Pro físico --- paridad completa de bench-coverage (cada
  backend corre en ambas plataformas; iOS retiene el transporte self-pipe de
  Phase-0 --- ver §6.3 --- y las optimizaciones SPSC ring buffer /
  \texttt{io\_uring} son Android-only), incluyendo el port nativo iOS de
  ObjectBox 5.3 cableado en esta revisión (entidades, bindings generados con
  Sourcery, y una test suite XCTest corrible en simulador en
  \texttt{experiment/storage/ios/Tests/ObjectBoxTests.swift}).
\item
  \textbf{Comparación de superficie primitiva y complejidad de desarrollo}
  contra cinco alternativas embebidas (SQLite, LMDB, RocksDB, ObjectBox,
  InMemory) más una referencia de Valkey-sobre-TCP: diez primitivas nativas
  contra \(\leq 4\), con medidas concretas de LOC para la misma capa de estado
  del agente (§5.6).
\item
  \textbf{Demostración end-to-end} vía una ablación 2×2 RAG sobre 200 queries NQ
  (§5.9, Tabla 15). Agregar la primitiva vectorial de Dazzle eleva
  \texttt{EM\_contains} 6.0× sobre Qwen 2.5 0.5B (0.105 → 0.630) y 6.7× sobre
  Qwen 2.5 1.5B (0.110 → 0.735); la celda de small-model-with-retrieval supera a
  la de large-model-without-retrieval en cada métrica factual. El resultado
  habla de qué se vuelve implementable cuando la primitiva correcta está
  presente, no de un delta en microsegundos sobre un competidor.
\item
  \textbf{Distribución pública del SDK} de \texttt{dazzle-sdk} 1.0.0-beta.4
  sobre cuatro registries estándar (Maven Central, pub.dev, npm, Swift Package
  Manager), con cinco adaptadores LLM detrás de una interfaz \texttt{LLMClient}
  común (Apéndice A). Apache-2.0 (Dazzle) / BSD-3-Clause (porciones derivadas de
  Valkey).
\end{enumerate}

\section{Contexto y motivación}\label{contexto-y-motivaciuxf3n}

\subsection{Requerimientos de estado para un agente LLM
on-device}\label{requerimientos-de-estado-para-un-agente-llm-on-device}

Un LLM-como-agente difiere de la inferencia de un solo turno en que debe:

\begin{itemize}
\tightlist
\item
  Mantener una ventana de observación bounded (buffer circular con desalojo
  automático).
\item
  Trackear running aggregates (min/max/avg sin re-scan O(N)).
\item
  Contar aproximadamente eventos distintos (HyperLogLog en 12 KB con
  \(\sim 0.8\) \% error relativo {[}@hyperloglog{]} vs O(N) exacto).
\item
  Aplicar TTL por campo a observaciones obsoletas (expiración automática sin un
  cron en la app).
\item
  Señalizar entre componentes (Pub/Sub sin polling).
\end{itemize}

Estas primitivas existen desde hace una década en Redis/Valkey. Lo nuevo es
embeberlas \textbf{dentro del proceso de la app móvil} para que el salto de red
entre agente y store desaparezca.

\subsection{Primitivas nativas por
backend}\label{primitivas-nativas-por-backend}

\textbf{Tabla 1 --- Disponibilidad por primitiva por backend embebido.} Cada
celda registra \emph{cómo} se alcanza la primitiva, no solamente \emph{si} es
alcanzable. \textbf{N} = llamada nativa de API tipada que el motor mismo provee.
\textbf{E} = extensión oficial provista por el proyecto upstream (e.g.,
R\emph{Tree de SQLite, en la amalgamation desde 2008). \textbf{A} =
re-implementación a nivel de aplicación --- el desarrollador escribe código
pegamento sobre primitivas más bajas en el motor. \textbf{T} = extensión de
tercero distribuida separadamente del motor. La fila ``Total Native'' cuenta
solo \textbf{N}; \textbf{E}, \textbf{A} y \textbf{T} significan que la primitiva
es }alcanzable* pero con costo de ingeniería distinto, capturado
cuantitativamente por el diferencial de LOC en la Tabla 9 y discutido
explícitamente en §6.3.

{\def\LTcaptype{none} % do not increment counter
\begin{longtable}[]{@{}
  >{\raggedright\arraybackslash}p{(\linewidth - 10\tabcolsep) * \real{0.4400}}
  >{\centering\arraybackslash}p{(\linewidth - 10\tabcolsep) * \real{0.1067}}
  >{\centering\arraybackslash}p{(\linewidth - 10\tabcolsep) * \real{0.1067}}
  >{\centering\arraybackslash}p{(\linewidth - 10\tabcolsep) * \real{0.1200}}
  >{\centering\arraybackslash}p{(\linewidth - 10\tabcolsep) * \real{0.1467}}
  >{\centering\arraybackslash}p{(\linewidth - 10\tabcolsep) * \real{0.0800}}@{}}
\toprule\noalign{}
\begin{minipage}[b]{\linewidth}\raggedright
Primitiva
\end{minipage} & \begin{minipage}[b]{\linewidth}\centering
Dazzle
\end{minipage} & \begin{minipage}[b]{\linewidth}\centering
SQLite
\end{minipage} & \begin{minipage}[b]{\linewidth}\centering
RocksDB
\end{minipage} & \begin{minipage}[b]{\linewidth}\centering
ObjectBox
\end{minipage} & \begin{minipage}[b]{\linewidth}\centering
LMDB
\end{minipage} \\
\midrule\noalign{}
\endhead
\bottomrule\noalign{}
\endlastfoot
Bounded stream con MAXLEN & N & A & A & A & A \\
Sorted sets con range query & N & A¹ & A & N & A \\
Atomic float increment & N & A² & A & A & A \\
TTL por campo & N & A & A & A & A \\
HyperLogLog & N & A & A & A & A \\
Geo indexing & N & E³ & A & A & A \\
Pub/Sub & N & A & A & N⁴ & A \\
Server-side scripting (Lua) & N & A & A & A & A \\
Vector search (HNSW) & N & T⁵ & A & N & A \\
Iteración basada en cursor & N & N & N & N & N \\
\textbf{Total Native (N)} & \textbf{10} & \textbf{1} & \textbf{1} & \textbf{4} &
\textbf{1} \\
\end{longtable}
}

¹ Implementable vía \texttt{ORDER\ BY\ …\ WHERE\ …\ BETWEEN} en SQL de la app. ²
Implementable vía \texttt{UPDATE\ …\ SET\ v\ =\ v\ +\ ?} dentro de una
transacción. ³ R*Tree es una extensión oficial de SQLite incluida dentro de la
amalgamation desde 2008. ⁴ El \texttt{DataSubscription} de ObjectBox es una API
de 1-call tipada para consumidores de notificación de cambios; lo contamos como
un equivalente nativo de pub/sub para el caso de uso del bucle del agente (un
solo consumidor suscrito a un stream tipado de cambios), aun cuando es
observer-pattern en lugar de topic-routed. ⁵ SQLite no tiene una primitiva
vectorial nativa ni oficial; el vector search lo proveen solo extensiones de
tercero (\texttt{sqliteai/sqlite-vector} comercial, \texttt{asg017/sqlite-vec}
open-source), que benchmarkearemos separadamente en §5.8 como peers comerciales
en lugar de como feature nativo de SQLite.

El conteo del headline (10 vs \(\leq 4\) nativas) mide \textbf{cuántas
primitivas el autor del agente puede alcanzar con una sola llamada tipada de
librería} --- el eje de ``superficie primitiva'' de la tesis Dazzle. Las
primitivas marcadas \textbf{E}, \textbf{A} o \textbf{T} en backends no-Dazzle
son \emph{implementables}; la pregunta es el costo de ingeniería, que medimos
directamente en §5.6 (Tabla 9: 175 LOC Android / 186 LOC iOS para Dazzle vs
200--290 LOC para los equivalentes basados en SQLite). Un revisor leyendo la
Tabla 1 como ``primitivas que existen en cualquier parte del ecosistema de este
motor'' llega a un gap mucho más pequeño. Esa es una pregunta distinta (y más
débil) que la que la Tabla 1 está estructurada para responder, y §6.3
Limitaciones deja la elección del eje explícita. Los experimentos de performance
que siguen respaldan el argumento cuantitativo, incluyendo el benchmark
vectorial head-to-head contra ObjectBox 4.x y SQLiteAI sqlite-vector 0.9.95 en
§5.8.

\subsection{Por qué embeber Valkey, y por qué como
fork}\label{por-quuxe9-embeber-valkey-y-por-quuxe9-como-fork}

Valkey fue diseñado para servidores Linux: daemon long-lived, TCP loopback,
asunciones POSIX de GNU libc, shutdown vía \texttt{exit()}. Ninguna sobrevive un
port a un proceso de app móvil. Dazzle aplica tres cambios como overlay del
fork. El primero es \textbf{portabilidad}: shims para Bionic libc de Android y
para el SDK de Apple. El segundo es \textbf{lifecycle}: el servidor no puede
matar a su proceso host. El tercero es \textbf{transporte in-process}: el TCP
loopback es puro overhead cuando cliente y servidor comparten el address space.

El upstream de Valkey no se vendoriza. Lo descargamos en un tag fijo y le
aplicamos tres diffs (\textasciitilde180 líneas en total) antes de compilar.
Publicar solo los diffs preserva la licencia BSD-3-Clause de Valkey; el código
original de Dazzle (en \texttt{core/} y \texttt{sdk/}) es Apache-2.0.

\section{El sistema Dazzle}\label{el-sistema-dazzle}

\subsection{Arquitectura}\label{arquitectura}

Dazzle se organiza en tres capas:

\begin{itemize}
\tightlist
\item
  \textbf{\texttt{core/}} (Apache-2.0). Transporte in-process
  (\texttt{core/transport/dazzle\_transport.c}) + snapshot cache + worker pool
  (Android). C puro, compilado con el Valkey patcheado.
\item
  \textbf{\texttt{versions/\textless{}ver\textgreater{}/patches/}} (diffs de
  build-time). Tres patches: Android (shims de Bionic), iOS (ajustes para el SDK
  de Apple + un \texttt{malloc\_zone\_t} custom para reportar RSS aislado), y un
  hook de 14 líneas en \texttt{server.c} que inserta
  \texttt{dazzle\_direct\_init()} después de \texttt{InitServerLast()}.
\item
  \textbf{\texttt{sdk/android/}, \texttt{sdk/ios/}} (Apache-2.0). Bridges JNI +
  Kotlin (AAR) y bridges C + Swift (XCFramework).
\end{itemize}

\subsection{Transport path}\label{transport-path}

El event loop de Valkey corre sobre un \texttt{pthread} dedicado. El thread de
la aplicación se comunica vía:

\begin{itemize}
\tightlist
\item
  \textbf{Pipe in-process.} El thread de la app hace \texttt{write()} de un
  pointer de 8 bytes a un \texttt{DirectRequest}, muy por debajo de
  \texttt{PIPE\_BUF} (\(\geq\) 512 bytes por POSIX, 4 096 en Linux), así que el
  kernel garantiza entrega atómica sin lock del lado escritor. El event loop
  despierta en \texttt{read()}, ejecuta \texttt{call()} con un cliente fake, y
  señaliza el \texttt{condvar} con la respuesta RESP.
\item
  \textbf{Snapshot cache write-through.} Cada comando mutador que cruza el pipe
  actualiza un cache in-process bajo rwlock antes de liberar al caller. Las
  lecturas subsiguientes adquieren el rdlock y retornan valores directamente,
  sin wake-up del event loop. El pipe es síncrono, así que para cuando una
  lectura se dispara el cache ya refleja cada escritura previa: \textbf{cero
  staleness por construcción}.
\item
  \textbf{Worker pool (Android).} 2--4 threads adicionales absorben lecturas
  concurrentes bajo rwlocks striped por slot cuando el flag de entorno
  \texttt{DAZZLE\_PARALLEL\_READS=1} está activo.
\end{itemize}

\subsection{Auto-mirror post-EVAL}\label{auto-mirror-post-eval}

Los scripts Lua (\texttt{EVALSHA}) ejecutan sus escrituras internas a través del
\texttt{call()} de Valkey, que no dispara el hook del snapshot-mirror. Sin
trabajo adicional, los backends que escriben campos dentro de un script Lua y
los leen después pagarían el costo pipe-HMGET en cada lectura. Dazzle hidrata el
cache automáticamente después de cada \texttt{EVAL}: itera sobre los KEYS
declarados del script, re-lee cada hash a través de la API de kvstore, y hace
upsert de sus campos al bucket correspondiente. Un backend que escribe 8 campos
dentro de un Lua EVALSHA los ve reflejados en el snapshot en la siguiente
lectura, sin trabajo adicional en el SDK.

\section{Diseño experimental}\label{diseuxf1o-experimental}

\subsection{Carga de trabajo}\label{carga-de-trabajo}

Agente sintético de monitoreo IoT: 200 lecturas con cinco campos (temperatura,
humedad, presión, batería, vibración), 11 anomalías inyectadas, 10 checkpoints
de inferencia espaciados uniformemente. Dos condiciones: \textbf{Stateless} (el
modelo recibe solo las últimas 20 lecturas) y \textbf{Augmented} (las 20
lecturas + un bloque de contexto materializado con agregados sobre la ventana
completa).

\subsection{Backends}\label{backends}

Cinco alternativas embebidas (SQLite, LMDB, ObjectBox, RocksDB, InMemory) más
una referencia Valkey-sobre-RESP-TCP (el servidor Valkey 8 sin modificaciones
con loopback localhost, incluido como baseline de ``transport overhead'') y
siete variantes de Dazzle exponiendo distintos paths de primitivas: básico, Lua,
Pipeline, HFE, HLL, \textbf{Precompute} (agregados materializados) y
\textbf{Vector} (búsqueda vectorial HNSW). Doce backends en total por
dispositivo, con \textbf{paridad completa de bench-coverage} sobre los sweeps
storage-only y vector-bench en esta revisión (cada backend corre en ambas
plataformas; el port iOS de ObjectBox 5.3 aterriza en §5.6, y el path
Dazzle-Vector corre en ambas plataformas vía el preset unificado
\texttt{paper384\_scale}).

\subsection{Dispositivos}\label{dispositivos}

Mediciones físicas en dos dispositivos:

\begin{itemize}
\tightlist
\item
  \textbf{Moto G35 5G} (Unisoc T760, 8×ARM @ 2.0/2.2 GHz, 4 GB LPDDR4X, Android
  14). El handset Android de presupuesto / referencia.
\item
  \textbf{iPhone 12 Pro} (Apple A14 Bionic, 6 GB LPDDR4X, iOS 26.3). El punto de
  validación cross-platform.
\end{itemize}

Los runs storage-only se hacen sobre ambos dispositivos. Los runs vector-bench
se hacen sobre el Moto G35 5G y, en esta revisión, también sobre el iPhone 12
Pro vía el preset \texttt{paper384\_scale}. El RAG end-to-end de §5.9 cubre dos
SoCs Android físicos (Moto G35 5G + Moto G30; §5.9.5) y aún no iPhone 12 Pro.

\subsection{Métricas}\label{muxe9tricas}

\begin{itemize}
\tightlist
\item
  \textbf{Ingest}: microsegundos por lectura (per-reading ingest µs).
\item
  \textbf{Retrieval}: microsegundos promedio sobre la ventana de
  retrieval-sample, p50 y p95.
\item
  \textbf{Footprint}: bytes en disco después de los 200 ingests
  (\texttt{stat.st\_blocks\ ×\ 512}, equivale a \texttt{du\ -k}).
\item
  \textbf{Vector}: p50 / p95 / p99 search latency sobre 100 queries por celda en
  estado estable warm; recall@k contra una verdad brute-force SQLite.
\item
  \textbf{RAG E2E}: \texttt{EM\_short}, \texttt{EM\_contains},
  \texttt{F1\_short}, \texttt{F1\_passage} por query, agregados sobre 200
  queries de NQ.
\end{itemize}

Bootstrap percentile-method (\(B = 10\,000\), seed = 42) {[}@efronbootstrap;
@davisonhinkley{]} sobre las arrays per-query para los CIs reportados en las
tablas. Los scripts de re-derivación están en
\texttt{research/scripts/bootstrap\_*\_lats.py}.

\section{Evaluación}\label{evaluaciuxf3n}

\subsection{La memoria externa habilita la
síntesis}\label{la-memoria-externa-habilita-la-suxedntesis}

\textbf{Tabla 2 --- Accuracy de síntesis factual sobre estadísticas globales
(checkpoint final).}

{\def\LTcaptype{none} % do not increment counter
\begin{longtable}[]{@{}
  >{\raggedright\arraybackslash}p{(\linewidth - 8\tabcolsep) * \real{0.2703}}
  >{\centering\arraybackslash}p{(\linewidth - 8\tabcolsep) * \real{0.2342}}
  >{\centering\arraybackslash}p{(\linewidth - 8\tabcolsep) * \real{0.1892}}
  >{\centering\arraybackslash}p{(\linewidth - 8\tabcolsep) * \real{0.2252}}
  >{\centering\arraybackslash}p{(\linewidth - 8\tabcolsep) * \real{0.0811}}@{}}
\toprule\noalign{}
\begin{minipage}[b]{\linewidth}\raggedright
Condición
\end{minipage} & \begin{minipage}[b]{\linewidth}\centering
Total de lecturas
\end{minipage} & \begin{minipage}[b]{\linewidth}\centering
Conteo de anomalías
\end{minipage} & \begin{minipage}[b]{\linewidth}\centering
Temperatura media
\end{minipage} & \begin{minipage}[b]{\linewidth}\centering
Score
\end{minipage} \\
\midrule\noalign{}
\endhead
\bottomrule\noalign{}
\endlastfoot
Stateless & -- (confabula \textasciitilde50) & -- (confabula 2--3) & --
(confabula \textasciitilde21 °C) & \textbf{0/3} \\
Augmented (cualquier backend) & Y (200) & Y (11) & Y (22.4 °C) & \textbf{3/3} \\
\end{longtable}
}

Este es el hallazgo empírico central del paper. Sin memoria externa el modelo no
tiene acceso a datos fuera de su ventana de contexto actual (las últimas 20
lecturas), así que las ``estadísticas globales'' que produce son
confabulaciones. Con un bloque de contexto materializado, el modelo lee los
valores exactos y los reproduce. El backend específico es irrelevante para este
hallazgo. Lo que importa es que la memoria exista y que el contexto sea preciso.

La implicación para el diseño de agentes es directa: \textbf{la capa de memoria
externa es infraestructura obligatoria para cualquier tarea que razone sobre
estado acumulado}, no una optimización opcional. La elección del backend después
se reduce a performance, superficie primitiva, y ergonomía.

\subsection{Caracterización del envelope de
latencia}\label{caracterizaciuxf3n-del-envelope-de-latencia}

Esta sección establece el envelope de latencia en el que un backend de agente
on-device debe operar. No es una declaración de ganador de microbenchmark.
Medimos cada backend en dos regímenes: un sweep storage-only en N = 200 lecturas
en ambos dispositivos físicos (Tabla 3), y una referencia bajo carga concurrente
en N = 20 000 (1 escritor + 1 lector, Moto G35 5G) llevada como un stress test
de estado estable. Los dos juntos cubren el envelope operativo realista de un
agente on-device long-running.

\textbf{Nota metodológica para las Tablas 3--6.} Las filas de la familia SQLite
(\texttt{sqlite}, \texttt{sqlite-optimized}, \texttt{sqlite-precompute}) vienen
de un sweep de mitigación aislado (Android: 3 rondas por backend, orden
randomizado; iOS: 3 rondas, dataset \texttt{dataset\_iot\_baseline}). Las filas
SQLite de la familia vector (\texttt{sqlite-vec} y variantes SQLiteAI) vienen
del harness vectorial alineado en dim = 384, k = 10, p50 search latency sobre
100 queries. El path SQLiteAI iOS requirió vendorizar la amalgamation de SQLite
con \texttt{SQLITE\_ENABLE\_LOAD\_EXTENSION=1} dentro del bench target, porque
el libsqlite3 del sistema de Apple viene con extension loading deshabilitado ---
esto está documentado como parte del artefacto (§8.1,
\texttt{experiment/backends/ios/sqlitevectorai/svai\_ios.c}) y se validó
end-to-end con la suite XCTest corrible en simulador antes del run en device.

\textbf{Tabla 3 --- Storage backends, envelope ingest/retrieval en N = 200 en
dispositivos físicos} (ingest µs por lectura / retrieval µs promedio).

{\def\LTcaptype{none} % do not increment counter
\begin{longtable}[]{@{}
  >{\raggedright\arraybackslash}p{(\linewidth - 4\tabcolsep) * \real{0.2826}}
  >{\raggedleft\arraybackslash}p{(\linewidth - 4\tabcolsep) * \real{0.3478}}
  >{\raggedleft\arraybackslash}p{(\linewidth - 4\tabcolsep) * \real{0.3696}}@{}}
\toprule\noalign{}
\begin{minipage}[b]{\linewidth}\raggedright
Backend
\end{minipage} & \begin{minipage}[b]{\linewidth}\raggedleft
Moto G35 5G ingest / retrieval
\end{minipage} & \begin{minipage}[b]{\linewidth}\raggedleft
iPhone 12 Pro ingest / retrieval
\end{minipage} \\
\midrule\noalign{}
\endhead
\bottomrule\noalign{}
\endlastfoot
Dazzle (default) & 2 889 / 2 132 & 185 / 411 \\
Dazzle-Lua & 1 345 / 1 824 & 181 / 251 \\
Dazzle-Pipeline & 1 331 / 1 959 & 143 / 444 \\
Dazzle-HFE & 2 312 / 1 893 & 177 / 389 \\
Dazzle-HLL & 2 844 / 2 491 & 185 / 454 \\
\textbf{Dazzle-Precompute} & \textbf{1 404 / 34.7} & \textbf{267 / 129} \\
Valkey 8 (RESP-TCP) & 3 754 / 3 320 & 432 / 1 474 \\
SQLite (default) & 251.95 / 1 061.3 & 61.94 / 268.7 \\
SQLite-Optimized & 59.09 / 899.4 & 23.64 / 199.4 \\
SQLite-Precompute & 207.77 / 75.0 & 317.10 / 27.6 \\
LMDB & 172 / 641 & 100 / 346 \\
RocksDB & 294 / 922 & 169 / 424 \\
ObjectBox \(\ddagger\) & 1 839 / 778 & 4 840 / 638 \\
In-memory (struct) & 9.97 / 282 & 5.35 / 204 \\
\end{longtable}
}

\(\ddagger\) El ingest de ObjectBox iOS bajó de 24 780 µs/r a 4 840 µs/r (5.1×
speedup) una vez que envolvimos el cuerpo de cada \texttt{ingest()} y
\texttt{storeCheckpointDecision()} en
\texttt{store.runInTransaction\ \{\ ...\ \}}. El port Kotlin corre
\textasciitilde10 calls de \texttt{box.put()} / \texttt{box.query()} por lectura
y depende del write-tx batching implícito por-call de ObjectBox-Java para
alcanzar 1 839 µs/r; ObjectBox-Swift abre una transacción de escritura nueva por
cada call (lock + journal append + fsync), así que las mismas 10 llamadas sin
coalescing pagan el overhead del lock 10×. El gap remanente de
\textasciitilde2.6× iPhone-vs-Moto es consistente con el allocator iOS y el
overhead FFI observado en otras partes de la tabla (ver la anomalía del
footprint Dazzle iOS discutida en §6.3); no es un bug del bench-harness.

\textbf{Tabla 4 --- Vector backends, envelope ingest/retrieval en N = 200 en
dispositivos físicos} (ingest µs por lectura / p50 search latency µs, dim = 384,
k = 10).

{\def\LTcaptype{none} % do not increment counter
\begin{longtable}[]{@{}
  >{\raggedright\arraybackslash}p{(\linewidth - 4\tabcolsep) * \real{0.3103}}
  >{\raggedleft\arraybackslash}p{(\linewidth - 4\tabcolsep) * \real{0.3333}}
  >{\raggedleft\arraybackslash}p{(\linewidth - 4\tabcolsep) * \real{0.3563}}@{}}
\toprule\noalign{}
\begin{minipage}[b]{\linewidth}\raggedright
Backend
\end{minipage} & \begin{minipage}[b]{\linewidth}\raggedleft
Moto G35 5G ingest / search
\end{minipage} & \begin{minipage}[b]{\linewidth}\raggedleft
iPhone 12 Pro ingest / search
\end{minipage} \\
\midrule\noalign{}
\endhead
\bottomrule\noalign{}
\endlastfoot
\textbf{Dazzle-Vector (HNSW)} & \textbf{74 / 65} & \textbf{25.17 / 18} \\
sqlite-vec default & 23.62 / 916.0 & 27.64 / 196 \\
sqlite-vec optimized & 9.66 / 860.0 & 20.67 / 193 \\
sqlite-vec precompute & 9.35 / 816.7 & 19.72 / 181 \\
SQLiteAI default & 6.08 / 135.3 & 6.90 / 29 \\
SQLiteAI optimized & 5.98 / 135.0 & 6.82 / 29 \\
SQLiteAI precompute & 5.17 / 111.3 & 7.27 / 23 \\
\end{longtable}
}

Dazzle-Vector en el Moto G35 5G ahora reporta \textbf{74 µs/r ingest, 65 µs p50
search a recall 1.000 (HNSW, ef = 10)} en N = 200. La celda previa
\texttt{—\ /\ —} en esta fila se debió a un bug del harness en
\texttt{SqliteBruteforceVector.create()} (la truth-source SQLite brute-force
nunca borraba su database entre runs porque el delete de archivo apuntaba al
directorio incorrecto, así que la truth de top-k por config se computaba contra
\textasciitilde20 000 docs obsoletos acumulados a través de bench launches
previos; solo la última config en cualquier sweep recuperaba fortuitamente
recall por encima del piso de 0.95). Ambos fixes aterrizaron en esta revisión:
el delete de la truth-source ahora resuelve a
\texttt{context.getDatabasePath(...)} y \texttt{VectorBenchmark} emite
\texttt{FLUSHALL} entre configs y entre variantes de motor Dazzle así que cada
medición es per-engine y per-config en aislamiento. El fix está en el commit
\texttt{1e3d5f5} de \texttt{experiment/backends/android/core/}
(\texttt{SqliteBruteforceVector.kt} + \texttt{VectorBenchmark.kt}), y el bench
post-fix está committeado en
\texttt{research/benchmarks/results/Moto\_G35\_5G/vector/vecbench\_moto\_g35\_5G\_1777369156656.json}
(SHA-256
\texttt{5a64d3692da166d96306d456697e43bb89c27c07b515e253346c5b33bc5c9b5b}). Las
Tablas 4, 5, 11 y el N-sweep de §5.8.4 / companion-report leen de este único
JSON, así que los números celda-por-celda y la narrativa de recall vienen de un
solo run en device, no de una mezcla de snapshots pre-fix y post-fix.

\textbf{Observación 1.} Dazzle-Precompute fija el piso alcanzable cuando los
agregados materializados se exponen como primitiva nativa. El retrieval promedia
34.7 µs en N = 200 en Moto G35 5G y 47 µs en N = 20 000 bajo carga concurrente.
El costo de retrieval es plano en N porque el bloque de contexto materializado
se reconstruye en el path de escritura y se lee con una copia de campo de tiempo
constante. En iPhone 12 Pro el retrieval aterriza en 129 µs. Estos números son
la cota inferior que un backend puede alcanzar cuando el patrón de lectura del
agente está compilado dentro del schema, no lo que cualquier backend ``podría''
alcanzar en principio.

\textbf{Observación 2.} Una concesión honesta que la tabla ahora hace explícita,
sobre las nueve variantes de la familia SQLite (default/optimised/precompute
sobre storage SQLite, sqlite-vec, y SQLiteAI): el retrieval Android se mueve de
1 061.3 µs (\texttt{sqlite} default) a 899.4 µs (\texttt{sqlite-optimized}) y a
75.0 µs (\texttt{sqlite-precompute}). Eso cierra la mayor parte del gap de
storage-latency y refuerza el framing del paper. El diferenciador no es un
límite intrínseco del algoritmo SQLite, sino cuánta ceremonia de materialización
y código de mantenimiento debe poseer la aplicación.

\textbf{La Observación 3 es la que realmente importa.} Incluso la latencia de
retrieval más alta del envelope (SQLite a 3 374 µs bajo carga concurrente en
Moto G35 5G) es aproximadamente el 0.0001 \% de un solo turno de inferencia
Gemma 4 E2B {[}@gemma4{]} (1--3 segundos). La diferencia perceptible para el
usuario entre 50 µs y 3 374 µs en la capa del agente es cero; ambos quedan
empequeñecidos por la llamada de inferencia que se sienta inmediatamente después
del retrieval. Las comparaciones en microsegundos no son el eje sobre el cual
elegir un backend. Los criterios dominantes son la superficie primitiva (Tabla
1, §2.2) y la complejidad de desarrollo requerida para entregar una capa de
estado funcional (§5.6).

\subsubsection{Por qué seguimos reportando los microsegundos
exactos}\label{por-quuxe9-seguimos-reportando-los-microsegundos-exactos}

Un revisor puede razonablemente preguntar: si cada retrieval está por debajo del
piso de ruido de inferencia y las comparaciones en microsegundos no son el eje
de elección, ¿por qué §5 gasta doce tablas enumerando microsegundos? Dos
razones, ambas independientes de la tesis del agente:

\textbf{(i) Consumidores no-LLM pueden compartir la capa de storage.} Aunque
este paper no los benchmarkea, el mismo path de código Dazzle es
estructuralmente compatible con dashboards de analítica on-device, pipelines de
pre-agregación de sensores, y capas de resolución de conflictos de sync ---
cargas donde las primitivas de retrieval (agregados materializados, bounded
streams, HyperLogLog, HNSW) son las mismas pero donde la latencia de retrieval
\emph{no} está dominada por un forward pass de LLM. Reportamos los números
absolutos en microsegundos para que un lector evaluando esas cargas pueda
extraer la fila relevante de la Tabla 3 / Tabla 5 / Tabla 11 directamente sin
re-correr el bench. No hacemos ningún claim competitivo sobre la posición de
Dazzle en esas cargas en este momento.

\textbf{(ii) Los números absolutos son un sanity check sobre la arquitectura.}
Un backend que se vaya a 100 ms+ en N = 200 (un orden de magnitud por encima de
lo que reportamos aquí) sugeriría una regresión digna de investigar más allá del
piso de ruido LLM. El patrón Write-Heavy / Read-Light está supuesto a mantener
el retrieval acotado; si una revisión futura rompe eso, la columna exacta de
microsegundos en T3 es donde la regresión se manifiesta. La misma lógica para la
ablación de §5.4 y la tabla de evolución de performance de §5.7 en el companion
report: esos números son cómo sabemos que el patrón arquitectónico se mantiene
intacto a través de revisiones, no cómo recomendamos elegir un backend.

Esta sección por tanto establece que todos los backends medidos se ubican dentro
del envelope aceptable para un agente LLM on-device. El resto de la evaluación
se lee en dos tracks: \textbf{(a)} para un lector centrado en el agent-loop, la
superficie primitiva (Tabla 1) y la complejidad de desarrollo (Tabla 9) impulsan
la selección de backend; \textbf{(b)} para un lector no-LLM, los números por
motor en microsegundos son el input directo. Ambas lecturas vienen de los mismos
datos; la diferencia es qué eje tratará el lector como el headline.

\subsection{Retrieval constante en N y footprint de
almacenamiento}\label{retrieval-constante-en-n-y-footprint-de-almacenamiento}

\textbf{Tabla 5 --- Promedio de latencia de retrieval a través de tamaños de
dataset, Moto G35 5G} (µs). Sweep en N \(\in\) \{200, 1 000, 5 000, 20 000\}.
Cada celda en la Tabla 5 y la Tabla 5b es la media aritmética por-run de la
ventana de retrieval-sample (\texttt{retrieval\_avg\_us} en cada
\texttt{scale\_\textless{}backend\textgreater{}\_*.json}); para backends con
múltiples sweep runs en el repo (típicamente 3--4), la celda es la mediana de
esas medias por-run. La columna correspondiente \texttt{retrieval\_p50\_us}
también está persistida en cada JSON para cualquier lector que prefiera la
mediana sobre la media, pero los valores de la tabla a lo largo de esta sección
son medias.

{\def\LTcaptype{none} % do not increment counter
\begin{longtable}[]{@{}lrrrr@{}}
\toprule\noalign{}
Backend & N = 200 & N = 1 k & N = 5 k & N = 20 k \\
\midrule\noalign{}
\endhead
\bottomrule\noalign{}
\endlastfoot
\textbf{Dazzle-Precompute} & \textbf{37} & \textbf{29} & \textbf{29} &
\textbf{34} \\
\textbf{Dazzle-Incremental} & 325 & 163 & 155 & 146 \\
InMemory & 663 & 314 & 273 & 317 \\
SQLite (default) & 1 061.3 & 739.1 & 717.4 & 721.6 \\
SQLite-Optimized & 899.4 & 508.1 & 514.5 & 505.4 \\
SQLite-Precompute & 75.0 & 59.8 & 62.1 & 60.6 \\
sqlite-vec default & 916.0 & 1 445.7 & 7 271.0 & 27 850.3 \\
sqlite-vec optimized & 860.0 & 1 434.7 & 7 145.3 & 28 574.7 \\
sqlite-vec precompute & 816.7 & 1 430.7 & 7 421.3 & 26 505.3 \\
SQLiteAI default & 135.3 & 303.0 & 1 475.0 & 9 812.3 \\
SQLiteAI optimized & 135.0 & 302.3 & 1 591.7 & 9 569.7 \\
SQLiteAI precompute & 111.3 & 233.3 & 889.0 & 3 071.7 \\
LMDB & 1 130 & 841 & 1 398 & 1 986 \\
\end{longtable}
}

\textbf{Tabla 5b --- Promedio de latencia de retrieval a través de tamaños de
dataset, iPhone 12 Pro} (µs). Misma métrica (media aritmética de la latencia
per-sample de retrieval, \texttt{retrieval\_avg\_us}) y misma grilla de sweep
que la Tabla 5.

{\def\LTcaptype{none} % do not increment counter
\begin{longtable}[]{@{}lrrrr@{}}
\toprule\noalign{}
Backend & N = 200 & N = 1 k & N = 5 k & N = 20 k \\
\midrule\noalign{}
\endhead
\bottomrule\noalign{}
\endlastfoot
\textbf{Dazzle-Precompute} & \textbf{25.4} & \textbf{17.8} & \textbf{17.8} &
\textbf{17.8} \\
\textbf{Dazzle-Incremental} & 139.2 & 63.7 & 63.3 & 63.2 \\
InMemory & 148.9 & 149.3 & 106.4 & 55.8 \\
SQLite (default) & 73.4 & 47.5 & 47.7 & 46.7 \\
SQLite-Optimized & 75.9 & 38.2 & 32.9 & 33.1 \\
SQLite-Precompute & 6.88 & 3.75 & 3.79 & 3.83 \\
sqlite-vec default & 196 & 548 & 3 578 & 15 038 \\
sqlite-vec optimized & 193 & 547 & 3 535 & 15 200 \\
sqlite-vec precompute & 181 & 548 & 3 535 & 15 145 \\
SQLiteAI default & 29 & 96 & 441 & 2 850 \\
SQLiteAI optimized & 29 & 96 & 440 & 2 842 \\
SQLiteAI precompute & 23 & 79 & 355 & 1 407 \\
LMDB & 167.4 & 67.2 & 46.9 & 85.8 \\
\end{longtable}
}

Las Tablas 5 y 5b deben leerse como un resultado de patrón de implementación, no
como una carrera de backends. La propiedad relevante es \textbf{retrieval
constante en N} cuando el sistema sirve agregados materializados en lugar de
reconstruir contexto desde historia cruda en query time. La fila
Dazzle-Precompute en la Tabla 5 es la mediana de las medias por-run de cuatro
sweep runs. Dazzle-Precompute y Dazzle-Incremental proveen este patrón
nativamente: el trabajo de agregación se paga en escritura, y el retrieval lee
estado materializado acotado. SQLite también puede realizar retrieval
cuasi-constante cuando el mismo patrón se implementa explícitamente. Por tanto,
el comportamiento constante en N es una propiedad del approach de agregados
materializados, no una propiedad propietaria de un solo motor.

\subsubsection{Footprint de almacenamiento después de 200
ingests}\label{footprint-de-almacenamiento-despuuxe9s-de-200-ingests}

\textbf{Tabla 6 --- Footprint en disco después de 200 ingests, por dispositivo}
(KB).

{\def\LTcaptype{none} % do not increment counter
\begin{longtable}[]{@{}lrr@{}}
\toprule\noalign{}
Backend & Moto G35 5G (KB) & iPhone 12 Pro (KB) \\
\midrule\noalign{}
\endhead
\bottomrule\noalign{}
\endlastfoot
\textbf{Dazzle (default)} & \textbf{6.7} & \textbf{153.3} \\
Dazzle-Lua & 6.6 & 153.3 \\
Dazzle-Pipeline & 6.7 & 153.3 \\
Dazzle-HFE & 6.6 & 153.3 \\
Dazzle-HLL & 7.0 & 177.6 \\
Dazzle-Precompute & 7.7 & 153.5 \\
Valkey 8 (RESP-TCP) & 7.9 & 177.4 \\
LMDB & 72.0 & 176.0 \\
ObjectBox & 96.0 & 268.0 \\
RocksDB & 4 515.8 & 104.0 \\
SQLite & 374.2 & 3 635.2 \\
SQLite-Optimized & 494.6 & 3 391.5 \\
SQLite-Precompute & 495.6 & 176.6 \\
sqlite-vec default & 4.1 & 4.0 \\
sqlite-vec optimized & 4.1 & 4.0 \\
sqlite-vec precompute & 4.1 & 4.0 \\
SQLiteAI default & 12.3 & 12.0 \\
SQLiteAI optimized & 12.3 & 12.0 \\
SQLiteAI precompute & 12.3 & 12.0 \\
\end{longtable}
}

El ratio de headline es \textbf{Dazzle vs SQLite}:

\begin{itemize}
\tightlist
\item
  Moto G35 5G: 6.7 KB vs 374.2 KB --- \textbf{56× más pequeño}
\item
  iPhone 12 Pro: 153.3 KB vs 3 635.2 KB --- \textbf{24× más pequeño}
\end{itemize}

El footprint es la restricción binding sobre hardware IoT y móvil de presupuesto
donde el budget de flash y el bandwidth de DRAM son ambos escasos. Un payload
\textbf{24--56× más pequeño} es la diferencia entre que los datos de Dazzle
quepan en un puñado de líneas de caché L2 por retrieval, versus barrer múltiples
páginas de SQLite por query --- lo que se compone hacia atrás en la imagen de
retrieval-latency de §5.2.

\subsection{Ablación factorial 2³: ¿qué capa
contribuye?}\label{ablaciuxf3n-factorial-2uxb3-quuxe9-capa-contribuye}

El stack Dazzle compone tres optimizaciones independientes: snapshot cache
(on/off), bucket dispatch por hash (1 bucket vs 16), worker pool paralelo
(off/on). Medimos seis combinaciones válidas sobre la misma carga: K = 8
agentes, mix 80/20 read/write, 15 s por celda.

\textbf{Tabla 7 --- Throughput en K = 8, Moto G35 5G} (ops/s, mix 80/20
read/write, 15 s por celda).

{\def\LTcaptype{none} % do not increment counter
\begin{longtable}[]{@{}
  >{\raggedright\arraybackslash}p{(\linewidth - 12\tabcolsep) * \real{0.2651}}
  >{\raggedleft\arraybackslash}p{(\linewidth - 12\tabcolsep) * \real{0.1205}}
  >{\raggedleft\arraybackslash}p{(\linewidth - 12\tabcolsep) * \real{0.1084}}
  >{\raggedleft\arraybackslash}p{(\linewidth - 12\tabcolsep) * \real{0.1446}}
  >{\raggedleft\arraybackslash}p{(\linewidth - 12\tabcolsep) * \real{0.1205}}
  >{\raggedleft\arraybackslash}p{(\linewidth - 12\tabcolsep) * \real{0.1205}}
  >{\raggedleft\arraybackslash}p{(\linewidth - 12\tabcolsep) * \real{0.1205}}@{}}
\toprule\noalign{}
\begin{minipage}[b]{\linewidth}\raggedright
Backend
\end{minipage} & \begin{minipage}[b]{\linewidth}\raggedleft
baseline
\end{minipage} & \begin{minipage}[b]{\linewidth}\raggedleft
workers
\end{minipage} & \begin{minipage}[b]{\linewidth}\raggedleft
snap-ser
\end{minipage} & \begin{minipage}[b]{\linewidth}\raggedleft
snap-par
\end{minipage} & \begin{minipage}[b]{\linewidth}\raggedleft
hash-ser
\end{minipage} & \begin{minipage}[b]{\linewidth}\raggedleft
hash-par
\end{minipage} \\
\midrule\noalign{}
\endhead
\bottomrule\noalign{}
\endlastfoot
\texttt{dazzle-incremental} & 12 106 & 9 604 & \textbf{28 897} & 28 154 & 27 926
& 28 304 \\
\texttt{dazzle-precompute} & 19 676 & 10 988 & \textbf{38 830} & 38 754 & 38 367
& 38 156 \\
\end{longtable}
}

Tres observaciones cuantificadas:

\begin{enumerate}
\def\labelenumi{\arabic{enumi}.}
\tightlist
\item
  \textbf{El snapshot cache solo es la capa de throughput.}
  \texttt{snap-linear-serial} (snapshot on, 1 bucket, sin workers) es el pico en
  ambas filas: \textbf{38 830 ops/s} para precompute y \textbf{28 897 ops/s}
  para incremental. Capear el hash-index de 16 buckets y el worker pool paralelo
  encima del cache (la columna ``full-stack'' \texttt{hash-par}) realmente
  entrega \textbf{38 156 / 28 304 ops/s} --- \texttt{snap-ser} está 1.8 \% por
  encima de \texttt{hash-par} para precompute y 2.1 \% por encima para
  incremental. El stack completo pierde, no gana, contra \texttt{snap-ser}; los
  deltas están dentro de la varianza run-to-run, así que la lectura práctica es
  ``el snapshot cache por sí solo es suficiente y las otras dos capas no agregan
  nada medible sobre esta carga''.
\item
  \textbf{Los workers solo regresan performance.} Sin un cache contra el cual
  hacer dispatch, el worker pool es puro overhead de queueing: 44 \% más lento
  para precompute, 21 \% más lento para incremental.
\item
  \textbf{El hash-index es future-proofing.} El keyspace de 4--5 keys del
  benchmark cabe trivialmente en un solo bucket; el dispatch O(1) se vuelve
  visible solo cuando el keyspace crece a docenas de keys o cuando los nombres
  de campo comparten prefijos largos.
\end{enumerate}

En una línea: \textbf{el snapshot cache es la única capa que mueve throughput
sobre esta carga. El worker pool paralelo y el dispatch hash-index son no-ops
aquí y solo ganan su costo sobre cargas con un keyspace más grande o más
contendido.}

\subsection{Validación cross-platform: iPhone 12
Pro}\label{validaciuxf3n-cross-platform-iphone-12-pro}

El mismo stack Dazzle corrido en un iPhone 12 Pro físico (iOS 26.3, A14 Bionic)
y un Moto G35 5G físico (Android 14, Unisoc T760), 200 lecturas, un solo
escritor + 100 muestras de retrieval (mismo protocolo storage-only que §5.2).

\textbf{Tabla 8 --- Resumen Dazzle cross-platform, Moto G35 5G vs iPhone 12 Pro}
(N = 200, un solo escritor + 100 muestras de retrieval).

{\def\LTcaptype{none} % do not increment counter
\begin{longtable}[]{@{}lrr@{}}
\toprule\noalign{}
Métrica & Moto G35 5G & iPhone 12 Pro \\
\midrule\noalign{}
\endhead
\bottomrule\noalign{}
\endlastfoot
Dazzle-Precompute ingest (µs/r) & 1 404 & 267 \\
Dazzle-Precompute retrieval avg & \textbf{34.7 µs} & \textbf{129 µs} \\
Dazzle-Precompute retrieval P50 & 34.3 µs & 139 µs \\
Dazzle-Precompute retrieval P95 & 42.2 µs & 189 µs \\
Dazzle-default retrieval avg & 2 132 µs & 411 µs \\
Ganancia de retrieval (precompute/default) & 61× & 3.2× \\
Footprint (Dazzle default) & 6.7 KB & 153.3 KB \\
Footprint vs SQLite & 56× menor & 24× menor \\
\end{longtable}
}

Tres observaciones sobre el comportamiento cross-platform:

\begin{enumerate}
\def\labelenumi{\arabic{enumi}.}
\item
  \textbf{La forma arquitectónica generaliza.} Dazzle-Precompute le gana al path
  Dazzle default en ambas plataformas --- el trade-off Write-Heavy / Read-Light
  se mantiene en ambas direcciones del par SoC + OS. El snapshot cache paga en
  cada lectura sin importar dónde corra.
\item
  \textbf{Los números absolutos de retrieval son sensibles a la plataforma.}
  Moto G35 5G aterriza en 34 µs P50; iPhone 12 Pro en 139 µs P50. El gap de 4×
  lo impulsa el path RESP del Valkey-fork in-process en iOS (tokio runtime +
  bridge Rust→Swift por retrieval), no el snapshot cache en sí. El path Android
  usa un bridge JNI más delgado con menos hops por call.
\item
  \textbf{La dirección del ingest se invierte.} El Moto necesita 1 404 µs por
  lectura; el iPhone necesita 267 µs (5.3× más barato en iPhone). Los cores más
  rápidos del A14 Bionic más el bandwidth LPDDR4X tragan la secuencia
  per-reading de \texttt{XADD} + \texttt{HINCRBYFLOAT} + \texttt{HSET} más
  rápido que el tier de presupuesto Unisoc T760.
\end{enumerate}

El resumen: la \textbf{contribución arquitectónica generaliza} sobre hardware y
OS (el costo se desplaza de read a write, el snapshot cache absorbe N, el
footprint se mantiene a escala de un solo motor), mientras que \textbf{los
números absolutos varían por SoC y stack de runtime}. Incluso sobre el path de
retrieval iPhone más lento, el P95 se mantiene bajo 200 µs. Eso es tres a cuatro
órdenes de magnitud bajo el forward pass LLM on-device (\(>10^{6}\) µs en
cualquier dispositivo), que sigue siendo el único término que importa al nivel
del agent-loop.

\subsection{Complejidad de desarrollo}\label{complejidad-de-desarrollo}

\textbf{Tabla 9 --- Líneas de código por backend, desagregadas por plataforma}
(SLOC: blanks y líneas de puro comentario excluidas; los counts Kotlin usan el
archivo Android \texttt{*ContextManager.kt} en
\texttt{experiment/backends/android/} y los counts Swift usan el archivo iOS
\texttt{*ContextManager.swift} en \texttt{experiment/backends/ios/}. Menor es
mejor; cada fila implementa la misma superficie de capa de estado del agente).

{\def\LTcaptype{none} % do not increment counter
\begin{longtable}[]{@{}
  >{\raggedright\arraybackslash}p{(\linewidth - 8\tabcolsep) * \real{0.1364}}
  >{\raggedleft\arraybackslash}p{(\linewidth - 8\tabcolsep) * \real{0.2091}}
  >{\raggedleft\arraybackslash}p{(\linewidth - 8\tabcolsep) * \real{0.1636}}
  >{\raggedleft\arraybackslash}p{(\linewidth - 8\tabcolsep) * \real{0.1273}}
  >{\raggedright\arraybackslash}p{(\linewidth - 8\tabcolsep) * \real{0.3636}}@{}}
\toprule\noalign{}
\begin{minipage}[b]{\linewidth}\raggedright
Backend
\end{minipage} & \begin{minipage}[b]{\linewidth}\raggedleft
Android (Kotlin) SLOC
\end{minipage} & \begin{minipage}[b]{\linewidth}\raggedleft
iOS (Swift) SLOC
\end{minipage} & \begin{minipage}[b]{\linewidth}\raggedleft
Suma (ambas)
\end{minipage} & \begin{minipage}[b]{\linewidth}\raggedright
Estilo de API
\end{minipage} \\
\midrule\noalign{}
\endhead
\bottomrule\noalign{}
\endlastfoot
\textbf{Dazzle} & \textbf{175} & \textbf{186} & \textbf{361} & Comandos de
primitivas tipados \\
InMemory & 174 & 172 & 346 & Manipulación directa de colecciones \\
LMDB & 200 & 227 & 427 & API byte-buffer \\
RocksDB & 203 & 227 & 430 & API byte-buffer con options \\
ObjectBox & 233 & 245 & 478 & Anotaciones de schema + DSL \\
SQLite & 268 & 292 & 560 & SQL + cursor mapping \\
\end{longtable}
}

El desglose por dos plataformas importa porque el port Android y el port iOS de
un backend se trackean entre sí dentro de \textasciitilde30 SLOC para cada fila
excepto SQLite (+24 en iOS, mayormente marshalling de result-row con C-API que
la API \texttt{Cursor} de Android oculta). El ordering relativo es estable a
través de plataformas: Dazzle es el backend más bajo no-InMemory en ambas,
SQLite es el más alto, y el gap entre Dazzle y SQLite es \textasciitilde93 SLOC
en Android y \textasciitilde106 SLOC en iOS --- ligeramente más amplio en iOS.

Dazzle requiere el menor número de líneas porque su API tipada (\texttt{XADD},
\texttt{HMGET}, \texttt{PFADD}) expresa directamente las operaciones del dominio
del agente sin schema translation ni cursor parsing.

\subsection{Evolución de performance
(resumen)}\label{evoluciuxf3n-de-performance-resumen}

El stack actual es el resultado de cinco rediseños de transporte in-process. La
contribución individual más grande es el \textbf{auto-mirror post-EVAL (+189 \%
sobre throughput de retrieval incremental} en K = 8, Moto G35 5G, mix 80/20
read-write), que hidrata el snapshot cache desde escrituras Lua-internas sin
ceremonia manual. Agregar el snapshot index buckedeado, el worker pool paralelo,
y el HSET inline en el script Lua contribuye un \textasciitilde14 \% adicional
sobre precompute. La tabla completa versión-por-versión (baseline 9 684 ops/s →
final 28 057 ops/s en incremental, 34 356 → 38 156 en precompute) vive en §1 del
companion engineering report
(\texttt{research/paper/companion\_engineering\_report.md}). Cada optimización
fue verificada como una contribución independiente; el cambio se persiste en el
repo y en el JSON de benchmark archivado, no se pierde a través de revisiones.

\subsection{Primitiva de búsqueda vectorial: envelope y operating
points}\label{primitiva-de-buxfasqueda-vectorial-envelope-y-operating-points}

El resultado que importa en esta sección es un envelope claim, no un
speed-competition claim: \textbf{la búsqueda vectorial HNSW de Dazzle entrega
retrieval sub-milisegundo a recall \(\geq 0.95\) a lo largo del rango operativo
estándar de mobile RAG (\(\dim \leq 384\), \(N \leq 10\,000\))}. Eso pone a
Dazzle en la misma clase algorítmica que ObjectBox 4.x, el peer comercial
mobile-first con soporte HNSW nativo: ambos son HNSW {[}@hnsw{]} con búsqueda
esperada \(O(\log N)\).

SQLiteAI sqlite-vector v0.9.95, la extensión SQLite comercial
production-shipping al cierre de Q1 2026, expone un linear scan acelerado por
SIMD cuantizado (\texttt{vector\_quantize\_scan}) como su único query path
optimizado; HNSW no es user-accessible en esta versión. Incluir SQLiteAI es
informativo como referencia de lo que los desarrolladores móviles del ecosistema
SQLite encuentran hoy, pero esta comparación abarca dos clases algorítmicas
(\(O(\log N)\) vs \(O(N)\)); el gap resultante refleja esa asimetría, no la
calidad de ingeniería de ninguna implementación.

El contrato del harness para cada celda de recall floor / operating-point en
esta sección está documentado, auditado, y post-mortemeado en el Apéndice B
(\texttt{research/paper/appendix\_b\_harness\_postmortem.md}). Cada celda en las
Tablas 4, 11 y 12 viene del run post-fix archivado en
\texttt{research/benchmarks/results/Moto\_G35\_5G/vector/vecbench\_moto\_g35\_5G\_1777369156656.json}
(commit \texttt{1e3d5f5}, SHA-256 citada en §5.2).

\subsubsection{5.8.1 Setup --- dos presets de
harness}\label{setup-dos-presets-de-harness}

La evaluación vectorial corre el mismo harness (\texttt{VectorBenchmark.kt} en
Android, \texttt{VectorBenchmark.swift} en iOS) contra dos grillas de
configuración distintas. El resto de §5.8 cita cualquiera dependiendo de la
pregunta a responder:

{\def\LTcaptype{none} % do not increment counter
\begin{longtable}[]{@{}
  >{\raggedright\arraybackslash}p{(\linewidth - 6\tabcolsep) * \real{0.4110}}
  >{\raggedleft\arraybackslash}p{(\linewidth - 6\tabcolsep) * \real{0.1233}}
  >{\raggedright\arraybackslash}p{(\linewidth - 6\tabcolsep) * \real{0.2740}}
  >{\raggedright\arraybackslash}p{(\linewidth - 6\tabcolsep) * \real{0.1918}}@{}}
\toprule\noalign{}
\begin{minipage}[b]{\linewidth}\raggedright
Preset
\end{minipage} & \begin{minipage}[b]{\linewidth}\raggedleft
Configs
\end{minipage} & \begin{minipage}[b]{\linewidth}\raggedright
Usado en
\end{minipage} & \begin{minipage}[b]{\linewidth}\raggedright
Qué responde
\end{minipage} \\
\midrule\noalign{}
\endhead
\bottomrule\noalign{}
\endlastfoot
\texttt{DEFAULT\_CONFIGS} & 9 (3 dim \(\times\) 3 N) & §5.8.3 envelope claim &
``¿Dazzle se mantiene sub-milisegundo a lo largo del rango mobile RAG estándar
(dim \(\leq\) 384, N \(\leq\) 10 000)?'' \\
\texttt{vector-bench-paper384-scale} (Android) / \texttt{paper384\_scale} (iOS)
& 4 (1 dim \(\times\) 4 N) & T11 + T12 + T13 + T14 & ``¿Cuál es el operating
point en la escala headline del paper, y cómo escala a N = 20 000?'' \\
\end{longtable}
}

La grilla de 9 configs barre \textbf{\{16, 128, 384\}} dim \(\times\)
\textbf{\{500, 2 000, 10 000\}} N para caracterizar la forma del envelope. La
grilla de 4 configs fija dim = 384 y barre N \(\in\) \textbf{\{200, 1 000, 5
000, 20 000\}} así que el operating point reportado en T11 (N = 20 000) está en
la misma grilla que el N-sweep de la familia SQLite en T12--T14. Ambas grillas
corren k = 10 nearest neighbours, distancia coseno, recall floor recall@10
\(\geq 0.95\), p50 / p95 / p99 sobre 100 queries por celda en estado estable
warm (descartando cold cache).

Configuraciones:

\begin{itemize}
\tightlist
\item
  \textbf{Dazzle HNSW (float32)}: M=32, efConstruction=400, efRuntime variado en
  \{10, 50, 100, 200\}.
\item
  \textbf{Dazzle SQ8 (int8 + NEON SDOT)}: mismos build params HNSW, más
  cuantización scalar en el path de storage con NEON dot-product en el path de
  comparación.
\item
  \textbf{ObjectBox 4.x}: HNSW vía la anotación \texttt{@HnswIndex}, configurada
  con los mismos M=32, efConstruction=400, distancia coseno, vectores float32.
\item
  \textbf{SQLiteAI 0.9.95}: \texttt{vector\_init} con FLOAT32 + coseno,
  \texttt{vector\_quantize(max\_memory=50MB)} para construir el snapshot int8,
  \texttt{vector\_quantize\_preload} para warmear el cache, queries a través de
  \texttt{vector\_quantize\_scan} (linear scan sobre vectores cuantizados).
\end{itemize}

\subsubsection{5.8.2 Tabla de operating-point}\label{tabla-de-operating-point}

\textbf{Tabla 11 --- Comparación de operating point por motor vectorial en dim =
384, N = 20 000, k = 10, recall@10 floor \(\geq 0.95\)} (p50 search latency en
µs, ingest total en ms, DB size en disco para motores de la familia SQLite,
\texttt{INFO\ memory\ →\ used\_memory\_dataset} para las filas in-memory
Dazzle/Valkey).

{\def\LTcaptype{none} % do not increment counter
\begin{longtable}[]{@{}
  >{\raggedright\arraybackslash}p{(\linewidth - 12\tabcolsep) * \real{0.2131}}
  >{\raggedright\arraybackslash}p{(\linewidth - 12\tabcolsep) * \real{0.1885}}
  >{\raggedright\arraybackslash}p{(\linewidth - 12\tabcolsep) * \real{0.0902}}
  >{\raggedleft\arraybackslash}p{(\linewidth - 12\tabcolsep) * \real{0.0902}}
  >{\raggedleft\arraybackslash}p{(\linewidth - 12\tabcolsep) * \real{0.1148}}
  >{\raggedleft\arraybackslash}p{(\linewidth - 12\tabcolsep) * \real{0.1066}}
  >{\raggedleft\arraybackslash}p{(\linewidth - 12\tabcolsep) * \real{0.1967}}@{}}
\toprule\noalign{}
\begin{minipage}[b]{\linewidth}\raggedright
Engine / Variant
\end{minipage} & \begin{minipage}[b]{\linewidth}\raggedright
Algorithm
\end{minipage} & \begin{minipage}[b]{\linewidth}\raggedright
Precision
\end{minipage} & \begin{minipage}[b]{\linewidth}\raggedleft
Recall@10
\end{minipage} & \begin{minipage}[b]{\linewidth}\raggedleft
p50 search ‡
\end{minipage} & \begin{minipage}[b]{\linewidth}\raggedleft
Ingest (ms)
\end{minipage} & \begin{minipage}[b]{\linewidth}\raggedleft
RAM / DB size
\end{minipage} \\
\midrule\noalign{}
\endhead
\bottomrule\noalign{}
\endlastfoot
\textbf{Dazzle SQ8} & HNSW + int8 SDOT & int8 & 0.959 & \textbf{208 µs} {[}203,
212{]} & 16 614 & 9.77 MB \(\dagger\) \\
Dazzle SQ8+Rerank & HNSW + int8 + fp32 & int8/fp32 & 0.982 & 330 µs & 14 966 &
39.06 MB \(\dagger\) \\
Dazzle F16 & HNSW + fp16 & fp16 & 0.984 & 297 µs & 26 082 & 17.09 MB
\(\dagger\) \\
Dazzle HNSW & HNSW & float32 & 0.952 & 343 µs & 43 028 & 31.74 MB \(\dagger\) \\
ObjectBox 4.x & HNSW & float32 & 0.994 & 853 µs {[}853, 1 078{]} & 387 405 & not
reported by runner \\
sqlite\_plain & linear scan & float32 & 1.000 & 707 302 µs & 1 614 & 79.45 MB \\
sqlite\_vec\_default & linear scan & float32 & 1.000 & 27 391 µs & 1 316 & 31.00
MB \\
sqlite\_vec\_optimized & linear scan & float32 & 1.000 & 28 575 µs & 869 & 31.00
MB \\
sqlite\_vec\_precompute & linear scan & float32 & 1.000 & 26 505 µs & 870 &
31.00 MB \\
SQLiteAI default \(*\) & quantized linear scan & int8 & 0.987 & 3 087 µs & 723 &
46.65 MB \\
SQLiteAI optimized \(*\) & quantized linear scan & int8 & 0.987 & 2 842 µs & 685
& 46.65 MB \\
SQLiteAI precompute \(*\) & quantized linear scan & int8 & 0.987 & 1 407 µs &
708 & 46.65 MB \\
\end{longtable}
}

‡ El intervalo \texttt{{[}lo,\ hi{]}} entre corchetes al lado de las dos celdas
headline de la misma clase (Dazzle SQ8, ObjectBox 4.x --- ambos HNSW) es un
\textbf{CI 95 \% bootstrap percentile-method cross-run sobre el estadístico p50}
{[}@efronbootstrap; @davisonhinkley{]} computado sobre los cuatro runs
independientes post-fix Moto G35 5G archivados en
\texttt{research/benchmarks/results/Moto\_G35\_5G/vector/} (timestamps
2026-04-28T05:05Z / 05:21Z / 09:24Z / 09:39Z; B = 10 000, seed = 42; script
\texttt{research/scripts/bootstrap\_vecbench\_cross\_run.py}; full report en
\texttt{research/paper/vecbench\_cross\_run\_ci.md}).

\(*\) Las tres filas SQLiteAI son \textbf{single-shot del run cross-engine} (el
mismo paper384\_scale pass que produjo cada otra fila en esta tabla). El sweep
dedicado de la familia SQLite de 3 rondas en el mismo N = 20 000 reporta
\texttt{default} 9 812 ± 363 µs, \texttt{optimized} 9 570 ± 548 µs,
\texttt{precompute} 3 072 ± 4 µs (companion report §2, Tabla 2). Mantenemos el
número cross-engine aquí para que cada fila de la Tabla 11 esté en el mismo
contexto de medición.

\(\dagger\) Las filas Dazzle/Valkey corren in-process y el bench no ejerció un
path de persistencia en disco (BGSAVE estaba deshabilitado para mantener el
costo de I/O fuera de la medición). Para poner las filas Dazzle en la misma
columna que los números \texttt{db\_file\_bytes} de la familia SQLite,
reportamos el \textbf{footprint analítico de RAM/disk} en su lugar, computado
como \texttt{(N\ ×\ dim\ ×\ bytes\_per\_component)\ +\ (N\ ×\ M\ ×\ 4)} --- los
vectores más el grafo HNSW (listas de vecinos int32, \texttt{M\ =\ 32} por
nodo).

El head-to-head justo es Dazzle vs ObjectBox: ambos motores corren HNSW.

A recall floor \(\geq 0.95\) (el operating point común de mobile RAG), Dazzle
SQ8 alcanza 208 µs p50 search en N = 20 000 contra ObjectBox 4.x en 853 µs ---
4.1× más rápido, ambos sobre HNSW. A recall estricto \(\geq 0.99\) la imagen se
invierte en el eje de recall: ObjectBox mantiene un recall más alto (0.994 vs
0.959 para Dazzle SQ8 en \(ef = 10\)) y Dazzle SQ8 se ubica en un operating
point de recall menor en esta configuración. Eso es un trade-off esperado entre
la cuantización int8 que Dazzle SQ8 usa y la precisión float32 que ObjectBox
mantiene, no una contradicción.

SQLiteAI sigue siendo una referencia secundaria útil porque es el path SQLite
production en muchos stacks móviles. El ratio observado contra SQLiteAI en N =
20 000 es consistente con comportamiento \(O(\log N)\) vs \(O(N)\) y debe leerse
con esa nota algorítmica adjunta.

\subsubsection{5.8.3 Interpretación del
envelope}\label{interpretaciuxf3n-del-envelope}

A lo largo de la grilla de 9 configuraciones, Dazzle se mantiene en el envelope
de search sub-milisegundo para el rango móvil objetivo y ObjectBox permanece en
el mismo orden de magnitud bajo settings HNSW alineados. La conclusión de
ingeniería relevante es que ambos sistemas son elecciones HNSW-class viables
para RAG on-device, con distintos trade-offs de constante-factor según el
operating point y el target de recall.

Versiones futuras de SQLiteAI sqlite-vector se espera que expongan HNSW. En ese
punto la comparación significativa se desplazará de asimetría de clase
algorítmica a análisis de constante-factor entre dos implementaciones HNSW.

\subsubsection{5.8.4 Sweep de variantes de extensión SQLite ---
resumen}\label{sweep-de-variantes-de-extensiuxf3n-sqlite-resumen}

\begin{quote}
\textbf{Disclosure de clase algorítmica.} Esta subsección compara Dazzle SQ8
(HNSW, search \(O(\log N)\) esperada) contra SQLiteAI sqlite-vector v0.9.95
(linear scan cuantizado acelerado por SIMD, \(O(N)\)). Los dos sistemas están en
\textbf{clases algorítmicas distintas}, no en la misma clase con constantes
diferentes. Los ratios numéricos reportados abajo son por tanto
\emph{referencias de envelope} --- describen qué tan adentro o afuera del
envelope de search sub-milisegundo aterriza cada path en un operating point
mobile-RAG representativo --- y \textbf{no claims de calidad de motor} sobre
SQLiteAI como producto.
\end{quote}

Para cerrar el gap del ``single SQLite path'' corrimos un N-sweep focalizado en
Moto G35 5G sobre backends vectoriales solo de la familia SQLite ---
\texttt{sqlite\_plain}, \texttt{sqlite\_vec} y SQLiteAI en sus variantes
\texttt{default} / \texttt{optimized} / \texttt{precompute} --- en dim = 384, k
= 10, 100 queries, 3 rondas, N \(\in\) \{200, 1 k, 5 k, 10 k, 20 k\}. La fila de
operating-point N = 20 000 para cada variante está mergeada en la Tabla 11
arriba.

La lectura es \textbf{sobre membresía del envelope, no una declaración de
ganador}. En N = 20 000, dim = 384, k = 10, recall \(\geq 0.95\):

\begin{itemize}
\tightlist
\item
  \textbf{Dazzle SQ8} aterriza en 208 µs (Tabla 11) --- claramente dentro del
  envelope de search sub-milisegundo.
\item
  \texttt{sqlite\_vector\_ai\_precompute} aterriza en \textbf{1 407 µs} en el
  cross-engine single-shot (Tabla 11) y en \textbf{3 072 ± 4 µs} medio a través
  del sweep dedicado de 3 rondas --- fuera del envelope sub-milisegundo, en
  milisegundos de un solo dígito.
\item
  \texttt{sqlite\_vec} (linear scan brute-force) está consistentemente por
  encima de ambas lecturas SQLiteAI.
\end{itemize}

El factor 6.8×--14.8× entre Dazzle SQ8 y SQLiteAI precompute en este N es por
tanto \emph{evidencia de que el path HNSW se mantiene dentro del envelope
sub-milisegundo mientras que el path linear-scan cruza al rango de los
milisegundos} --- i.e., documenta un threshold de tamaño de corpus más allá del
cual la elección del motor empieza a importar al nivel del agent-loop.

\subsection{Aplicación end-to-end en Moto G35 5G: small-model + RAG
on-device}\label{aplicaciuxf3n-end-to-end-en-moto-g35-5g-small-model-rag-on-device}

\begin{quote}
\textbf{Alcance de esta sección.} La ablación RAG end-to-end reportada abajo se
corrió originalmente sobre un solo dispositivo físico (Motorola Moto G35 5G,
Unisoc T760, cluster A76 ARMv8.2). §5.9.5 extiende el run a un segundo SoC
Android físico (Motorola Moto G30, Qualcomm Snapdragon 662, cluster A73 ARMv8.0)
y reproduce la 2×2 con CIs 95 \% superpuestos en cada celda --- tres de las
cuatro celdas son bit-idénticas (Q4\_K\_M + greedy decoding + retrieval idéntico
determinizan el token stream) y la cuarta difiere por 0.005 en el punto
estimado. La cobertura end-to-end en iPhone 12 Pro queda como trabajo futuro; la
validación storage-path en iPhone 12 Pro ya aparece en §5.5.
\end{quote}

La tesis de esta sección es a nivel de aplicación: la disponibilidad de una
primitiva vectorial in-process habilita un patrón de agente cualitativamente
distinto --- small-model + RAG on-device {[}@lewisrag; @karpukhindpr{]} --- que
puede alcanzar accuracy factual mayor que un modelo más grande sin retrieval. La
pipeline completa corre on-device, sin cloud y sin servidor externo, incluido en
un teléfono Android de \$150. El corpus de 2 K passages que indexamos también
mantiene cada passage retrieved bien dentro de la ventana de contexto del
small-model, así que los artefactos de precisión-de-retrieval causados por
degradación de contexto-largo {[}@liulostmiddle{]} no están en el path entre el
índice y la respuesta.

\subsubsection{5.9.1 Setup}\label{setup}

Indexamos 2 000 passages de un mini-slice determinista de Natural Questions
{[}@nq2019{]} en un Motorola Moto G35 5G (Unisoc T760, 4 GB RAM, Android 14).
Los passages y los pares (question, gold-passage) vienen del split SBERT pair de
NQ {[}@sbertnq{]}; los aliases de short-answer canónicos se joinean desde
\texttt{nq\_open} {[}@nqopen; @leeorqa{]} por texto de pregunta lowercased.

\textbf{Precondición de aliases.} No cada NQ question tiene short-answer
canónico en \texttt{nq\_open} --- el join produce un campo
\texttt{short\_answers} solo para el subset de questions que los autores de
NQ-open anotaron. Pre-filtramos el pool de candidatos para mantener solo
questions con al menos un short-answer alias, después sampleamos 200 de esos
(así cada query en el bench tiene un set de ground-truth no-vacío para
\texttt{EM\_short} / \texttt{F1\_short}). Los 2 000 passages se construyen como
los 200 gold positives más 1 800 distractors random sacados del mismo pool, seed
42; los passages se truncan a 1 800 caracteres para caber en el contexto de 512
tokens de \texttt{bge-small-en-v1.5}. El slice es regenerable desde
\texttt{research/scripts/nq\_slice.py} y está fingerprinted por el prefijo
sha256 \texttt{63be4b8894c71ff3} (provenance completa en
\texttt{research/data/nq\_slice/README.md}).

Embeddings producidos con \texttt{bge-small-en-v1.5} {[}@bgesmall{]} (q4\_k\_m,
24 MB, dim 384, n\_ctx 512); índice Dazzle HNSW\_SQ8 con M=32, efC=400. Para
cada uno de los 200 queries hacemos retrieve de los k=5 passages más
cosine-similar con efRuntime=64 y los inyectamos en el prompt. Configuraciones:

\begin{itemize}
\tightlist
\item
  \textbf{Qwen 2.5 0.5B Instruct + Dazzle RAG.} GGUF q4\_k\_m, 380 MB en disco.
  n\_ctx=2 048. Prompt formado como
  \texttt{\textless{}system\textgreater{}…\ \textless{}retrieved\ passages\textgreater{}…\ \textless{}user\ question\textgreater{}}.
  Decoding greedy, max\_new\_tokens=64.
\item
  \textbf{Qwen 2.5 1.5B Instruct sin RAG.} GGUF q4\_k\_m, 940 MB en disco ---
  2.5× más grande, 3× más parámetros. n\_ctx=2 048. Prompt formado como
  \texttt{\textless{}system\textgreater{}\ \textless{}user\ question\textgreater{}}.
  Decoding greedy, max\_new\_tokens=64.
\end{itemize}

\textbf{Protocolo de evaluación.} Para cada query el modelo genera hasta 64
tokens. Extraemos el span de respuesta predicho \texttt{y\_pred} como el prefijo
de la generación hasta la primera newline o eco fresco \texttt{Question:} /
\texttt{Answer:}. Tanto \texttt{y\_pred} como el short-answer gold
\texttt{y\_gold} se normalizan siguiendo el protocolo SQuAD v1.1 {[}@squad{]}
(lowercasing, removal de los artículos \texttt{\{a,\ an,\ the\}} y puntuación,
collapsing de whitespace).

Reportamos tres métricas de short-answer más un backup de passage:

\begin{itemize}
\tightlist
\item
  \textbf{\texttt{EM\_short}} (estricto) --- \texttt{1} iff
  \texttt{norm(y\_pred)\ ==\ norm(y\_gold)} para cualquier alias, \texttt{0} en
  otro caso.
\item
  \textbf{\texttt{F1\_short}} --- F1 a nivel token (whitespace tokenisation)
  entre \texttt{norm(y\_pred)} y el alias \texttt{norm(y\_gold)} que mejor
  matchea, promediado sobre queries.
\item
  \textbf{\texttt{EM\_contains}} (laxo) --- \texttt{1} iff los tokens
  normalizados de cualquier alias aparecen como substring contiguo en alguna
  parte de la generación \emph{completa}. Esto separa \emph{``el modelo sabe la
  respuesta''} de \emph{``el modelo responde concisamente''}.
\item
  \textbf{\texttt{F1\_passage}} --- token-F1 entre \texttt{y\_pred} (whitespace
  tokens, sin normalización agresiva) y el passage gold completo; mide si la
  generación reproduce el contexto de soporte.
\end{itemize}

\subsubsection{5.9.2 Resultado --- matriz 2×2
completa}\label{resultado-matriz-22-completa}

\textbf{Tabla 15 --- Accuracy factual end-to-end RAG sobre 200 queries NQ} (Moto
G35 5G; decoding greedy, max\_new\_tokens = 64; mismo bench run para las cuatro
celdas; raw JSON en
\texttt{research/benchmarks/results/Moto\_G35\_5G/rag\_2x2/rag\_e2e\_moto\_g35\_5G\_1777395311213.json},
SHA-256
\texttt{00d21f6c8752ffaa1015624b69a5e5d0fd403670d72561e3838bdac0ab461e76}).

Cada celda se reporta como point estimate con el CI 95 \% bootstrap
percentile-method \texttt{{[}lo,\ hi{]}} sobre la array per-query de la métrica
(n = 200, B = 10 000, seed = 42). El report bootstrap completo (CIs por celda,
ratio CIs paired-qid sobre las 16 celdas métrica × ratio, flags de
significancia) vive en \texttt{research/paper/rag\_2x2\_with\_ci.md} y es
regenerable desde el mismo JSON vía
\texttt{research/scripts/bootstrap\_rag\_metrics.py}.

{\def\LTcaptype{none} % do not increment counter
\begin{longtable}[]{@{}
  >{\raggedright\arraybackslash}p{(\linewidth - 10\tabcolsep) * \real{0.1929}}
  >{\raggedright\arraybackslash}p{(\linewidth - 10\tabcolsep) * \real{0.0643}}
  >{\raggedright\arraybackslash}p{(\linewidth - 10\tabcolsep) * \real{0.1857}}
  >{\raggedright\arraybackslash}p{(\linewidth - 10\tabcolsep) * \real{0.1857}}
  >{\raggedright\arraybackslash}p{(\linewidth - 10\tabcolsep) * \real{0.1857}}
  >{\raggedright\arraybackslash}p{(\linewidth - 10\tabcolsep) * \real{0.1857}}@{}}
\toprule\noalign{}
\begin{minipage}[b]{\linewidth}\raggedright
Configuración
\end{minipage} & \begin{minipage}[b]{\linewidth}\raggedright
Tamaño
\end{minipage} & \begin{minipage}[b]{\linewidth}\raggedright
EM\_short
\end{minipage} & \begin{minipage}[b]{\linewidth}\raggedright
EM\_contains
\end{minipage} & \begin{minipage}[b]{\linewidth}\raggedright
F1\_short
\end{minipage} & \begin{minipage}[b]{\linewidth}\raggedright
F1\_passage
\end{minipage} \\
\midrule\noalign{}
\endhead
\bottomrule\noalign{}
\endlastfoot
Qwen 0.5B (sin RAG) & 380 MB & 0.015 {[}0.000, 0.035{]} & 0.105 {[}0.065,
0.150{]} & 0.079 {[}0.055, 0.106{]} & 0.151 {[}0.138, 0.164{]} \\
Qwen 0.5B + Dazzle RAG & 380 MB & \textbf{0.120} {[}0.080, 0.165{]} &
\textbf{0.630} {[}0.565, 0.695{]} & \textbf{0.235} {[}0.191, 0.283{]} & 0.334
{[}0.300, 0.369{]} \\
Qwen 1.5B (sin RAG) & 940 MB & 0.045 {[}0.020, 0.075{]} & 0.110 {[}0.070,
0.155{]} & 0.118 {[}0.084, 0.154{]} & 0.085 {[}0.073, 0.098{]} \\
Qwen 1.5B + Dazzle RAG & 940 MB & \textbf{0.220} {[}0.170, 0.275{]} &
\textbf{0.735} {[}0.670, 0.795{]} & \textbf{0.331} {[}0.282, 0.380{]} & 0.387
{[}0.355, 0.420{]} \\
\end{longtable}
}

Leyendo la matriz 2×2 a través de filas y columnas (los CIs de ratios son
bootstrap paired-qid; ★ marca aquellos cuyo CI 95 \% excluye 1.0):

\begin{itemize}
\tightlist
\item
  \textbf{Lift por adición de RAG} (mismo modelo, with-RAG vs no-RAG):
  \texttt{EM\_contains} 0.630 vs 0.105 en 0.5B (\textbf{6.0×} \(\bigstar\)),
  0.735 vs 0.110 en 1.5B (\textbf{6.7×} \(\bigstar\)); \texttt{EM\_short} 0.120
  vs 0.015 en 0.5B (\textbf{8.0×} \(\bigstar\)), 0.220 vs 0.045 en 1.5B
  (\textbf{4.9×} \(\bigstar\)).
\item
  \textbf{El modelo más pequeño con retrieval supera al modelo más grande sin
  retrieval en cada métrica factual.} 0.5B + RAG vs 1.5B sin RAG:
  \texttt{EM\_contains} 0.630 vs 0.110 (\textbf{5.7×}), \texttt{EM\_short} 0.120
  vs 0.045 (\textbf{2.7×}), \texttt{F1\_short} 0.235 vs 0.118 (\textbf{2.0×}),
  \texttt{F1\_passage} 0.334 vs 0.085 (\textbf{3.9×}).
\end{itemize}

\subsubsection{5.9.3 Latency vs accuracy
frontier}\label{latency-vs-accuracy-frontier}

La capacidad cruda del modelo y el retrieval contribuyen aditivamente. El 1.5B +
RAG (0.220 EM\_short / 0.735 EM\_contains) es el máximo global, pagando un costo
de p50 per-turn de 49 s. Para contexto del frontier:

\begin{itemize}
\tightlist
\item
  0.5B sin RAG: 2.50 s p50 per-turn, 0.015 EM\_short.
\item
  0.5B + RAG: 17.6 s p50 per-turn, 0.120 EM\_short.
\item
  1.5B sin RAG: 2.98 s p50 per-turn, 0.045 EM\_short.
\item
  1.5B + RAG: 49.0 s p50 per-turn, 0.220 EM\_short.
\end{itemize}

El gap dominante de latencia entre filas RAG y no-RAG son los \textasciitilde5
KB de passages retrieved (k=5, \textasciitilde1 KB cada uno) que el modelo tiene
que prefill en cada turno. En decoding token los runs RAG y no-RAG decodifican
aproximadamente el mismo número de tokens (\texttt{max\_new\_tokens=64}); el
delta de latencia es prácticamente todo prefill cost. El path de
retrieval-de-storage en sí (HNSW search sobre 2 000 passages) es
\textasciitilde0.6 ms en wall clock --- invisible en este gráfico.

Lo que importa es que el \textbf{patrón RAG está disponible on-device sin
dependencia de cloud}; la contribución del backend de storage a ese turn budget
es el costo de \textasciitilde22 ms + 0.6 ms embed-and-search --- bajo el 0.05
\% del row RAG total más rápido.

\subsubsection{5.9.5 Extensión cross-platform a tres SoCs Android y un SoC
Apple}\label{extensiuxf3n-cross-platform-a-tres-socs-android-y-un-soc-apple}

Para chequear que las conclusiones de §5.9 no son artefactos de un solo chip ---
ni de un solo OS --- re-corrimos la 2×2 completa sobre tres dispositivos
adicionales que cubren dos sistemas operativos:

\begin{itemize}
\tightlist
\item
  un Motorola Moto G30 (Qualcomm Snapdragon 662, cluster Cortex-A73, baseline
  ARMv8.0, 4 GB LPDDR4, Android 11);
\item
  un Huawei P20 Lite ANE-LX3 (HiSilicon Kirin 659, Cortex-A53, baseline ARMv8.0,
  4 GB LPDDR4, Android 9.1 / EMUI 9 / kernel 4.9);
\item
  un \textbf{Apple iPhone 12 Pro} (Apple A14 Bionic, cores Firestorm + Icestorm,
  ARMv8.4 + ISA Apple-private, 6 GB unified memory, iOS

  \begin{enumerate}
  \def\labelenumi{\arabic{enumi})}
  \setcounter{enumi}{25}
  \tightlist
  \item
    --- la única fila no-Android, ejercitando el port Swift de
    \texttt{RagE2EBench} en \texttt{experiment/llm/ios/} sobre los mismos entry
    points C \texttt{dazzle\_llama\_*} y \texttt{dazzle\_vs\_*} que el JNI
    Android usa en las otras tres filas.
  \end{enumerate}
\end{itemize}

Las cuatro filas corrieron con los mismos archivos de modelo, el mismo slice NQ
de 2 000 passages (prefijo sha256 \texttt{63be4b8894c71ff3}), y el mismo build
de Dazzle (\texttt{libdazzle.so} baseline en Android v8.0,
\texttt{libdazzle\_v82.so} en Android v8.2, static \texttt{libvalkey-server.a}
en iOS --- todos derivados del mismo árbol de fuentes
\texttt{core/platform/dazzle\_llama.cpp} +
\texttt{sdk/android/src/main/cpp/valkeysearch\_module.cc}). Greedy decoding,
passages retrieved idénticos, token streams idénticos.

\textbf{Tabla 17 --- Reproducción cross-platform §5.9. Medias \texttt{F1\_short}
con CIs 95 \% bootstrap-percentile (B = 10 000, seed = 42). Greedy decoding,
mismos pesos de modelo, mismo slice NQ. Las ratios con estrella colapsan el par
(numerador, denominador) en un único test de significancia contra 1.0 bajo
paired-qid resampling. Las filas del Kirin 659 e iPhone 12 Pro corren con
\texttt{Algorithm.FLAT} en lugar de \texttt{HNSW} (sidebar items 6 y 13); el
recall de FLAT a esta escala (N = 2 000, dim = 384, k = 5) es \textgreater{} 99
\%, por lo que \texttt{F1\_short} por celda difiere a lo más
\textasciitilde0.005 entre los dos algoritmos --- dentro del CI por celda.}

{\def\LTcaptype{none} % do not increment counter
\begin{longtable}[]{@{}
  >{\raggedright\arraybackslash}p{(\linewidth - 18\tabcolsep) * \real{0.0682}}
  >{\raggedright\arraybackslash}p{(\linewidth - 18\tabcolsep) * \real{0.0795}}
  >{\raggedright\arraybackslash}p{(\linewidth - 18\tabcolsep) * \real{0.0568}}
  >{\raggedright\arraybackslash}p{(\linewidth - 18\tabcolsep) * \real{0.0568}}
  >{\raggedright\arraybackslash}p{(\linewidth - 18\tabcolsep) * \real{0.0568}}
  >{\raggedright\arraybackslash}p{(\linewidth - 18\tabcolsep) * \real{0.0682}}
  >{\raggedright\arraybackslash}p{(\linewidth - 18\tabcolsep) * \real{0.1591}}
  >{\raggedright\arraybackslash}p{(\linewidth - 18\tabcolsep) * \real{0.1477}}
  >{\raggedright\arraybackslash}p{(\linewidth - 18\tabcolsep) * \real{0.1591}}
  >{\raggedright\arraybackslash}p{(\linewidth - 18\tabcolsep) * \real{0.1477}}@{}}
\toprule\noalign{}
\begin{minipage}[b]{\linewidth}\raggedright
Chip
\end{minipage} & \begin{minipage}[b]{\linewidth}\raggedright
µarch
\end{minipage} & \begin{minipage}[b]{\linewidth}\raggedright
ISA
\end{minipage} & \begin{minipage}[b]{\linewidth}\raggedright
RAM
\end{minipage} & \begin{minipage}[b]{\linewidth}\raggedright
OS
\end{minipage} & \begin{minipage}[b]{\linewidth}\raggedright
Algo
\end{minipage} & \begin{minipage}[b]{\linewidth}\raggedright
small no-RAG
\end{minipage} & \begin{minipage}[b]{\linewidth}\raggedright
small + RAG
\end{minipage} & \begin{minipage}[b]{\linewidth}\raggedright
large no-RAG
\end{minipage} & \begin{minipage}[b]{\linewidth}\raggedright
large + RAG
\end{minipage} \\
\midrule\noalign{}
\endhead
\bottomrule\noalign{}
\endlastfoot
Unisoc T760 (Moto G35 5G) & A76 & v8.2 & 6 GB & Android 14 & HNSW & 0.079
{[}0.055, 0.105{]} & 0.235 {[}0.191, 0.283{]} & 0.118 {[}0.084, 0.154{]} & 0.487
{[}0.431, 0.542{]} \\
QCOM SD662 (Moto G30) & A73 & v8.0 & 4 GB & Android 11 & HNSW & 0.079 {[}0.055,
0.105{]} & 0.240 {[}0.195, 0.287{]} & 0.118 {[}0.084, 0.154{]} & 0.487 {[}0.431,
0.542{]} \\
HiSi Kirin 659 (Huawei P20 Lite) & A53 & v8.0 & 4 GB & Android 9 & FLAT & 0.079
{[}0.055, 0.105{]} & 0.236 {[}0.191, 0.283{]} & 0.119 {[}0.085, 0.155{]} & 0.484
{[}0.427, 0.539{]} \\
\textbf{Apple A14 (iPhone 12 Pro)} & \textbf{Firestorm} & \textbf{v8.4+} &
\textbf{6 GB} & \textbf{iOS 26} & \textbf{FLAT} & \textbf{0.079 {[}0.055,
0.105{]}} & \textbf{0.236 {[}0.191, 0.283{]}} & \textbf{0.119 {[}0.085,
0.155{]}} & \textbf{0.488 {[}0.432, 0.543{]}} \\
\end{longtable}
}

Los tres chips reproducen las mismas cuatro celdas dentro del CI por celda. Tres
de las cuatro celdas (\texttt{small\ no-RAG}, \texttt{large\ no-RAG},
\texttt{large\ +\ RAG}) son bit-idénticas entre T760 y SD662; la fila del Kirin
matchea con ambos a ≤ 0.003 --- el delta de recall FLAT vs HNSW y la selección
de kernel ggml v8.0 vs v8.2 contribuyen juntos menos que el half-width del CI
con \texttt{B\ =\ 10\ 000} en cada celda.

La fila del Kirin 659 reporta el run tal como completó sobre el chip después de
una \textbf{investigación de quince passes} (deadlock de HNSW
\texttt{addPoint(label\ =\ 0)} → fallback a \texttt{Algorithm.FLAT}; acumulación
del kill-score de iAware en EMUI 9 a través del embed loop → driver
multi-process de cuatro fases; segfault SIMD del path FLAT del SDK sobre
Cortex-A53 → scan brute-force escalar portable sobre un mirror
\texttt{fp32\_store}; deadlock LLM en split-prefill sobre prompt de 570 tokens
con \texttt{n\_batch\ =\ 512} → \texttt{n\_batch\ =\ n\_ctx}; iAware mata el
bench process alrededor de la query 30/200 en las variantes Qwen 1.5B → 20
chunks de 10 queries cada uno, un \texttt{am\ instrument} fresh por chunk,
partials concatenados en el merge --- ver
\texttt{research/results/cross\_platform\_e2e/\ ane\_lx3\_kirin659\_investigation.md}).
Con los cinco fixes shippeados, la fila del Kirin es \textbf{directamente
comparable} a los chips HNSW y las ratios son idénticas:

{\def\LTcaptype{none} % do not increment counter
\begin{longtable}[]{@{}
  >{\raggedright\arraybackslash}p{(\linewidth - 6\tabcolsep) * \real{0.0667}}
  >{\raggedright\arraybackslash}p{(\linewidth - 6\tabcolsep) * \real{0.3111}}
  >{\raggedright\arraybackslash}p{(\linewidth - 6\tabcolsep) * \real{0.3111}}
  >{\raggedright\arraybackslash}p{(\linewidth - 6\tabcolsep) * \real{0.3111}}@{}}
\toprule\noalign{}
\begin{minipage}[b]{\linewidth}\raggedright
Chip
\end{minipage} & \begin{minipage}[b]{\linewidth}\raggedright
small + RAG / large no-RAG
\end{minipage} & \begin{minipage}[b]{\linewidth}\raggedright
small + RAG / small no-RAG
\end{minipage} & \begin{minipage}[b]{\linewidth}\raggedright
large + RAG / large no-RAG
\end{minipage} \\
\midrule\noalign{}
\endhead
\bottomrule\noalign{}
\endlastfoot
Unisoc T760 & 2.00× {[}1.44, 2.89{]} ★ & 2.97× {[}2.05, 4.51{]} ★ & 4.13×
{[}3.10, 5.86{]} ★ \\
QCOM SD662 & 2.03× {[}1.47, 2.95{]} ★ & 3.03× {[}2.09, 4.58{]} ★ & 4.13×
{[}3.10, 5.86{]} ★ \\
HiSi Kirin 659 & 1.98× {[}1.43, 2.88{]} ★ & 2.98× {[}2.06, 4.54{]} ★ & 4.07×
{[}3.04, 5.74{]} ★ \\
\textbf{Apple A14} & \textbf{1.98× {[}1.43, 2.88{]} ★} & \textbf{2.98× {[}2.06,
4.54{]} ★} & \textbf{4.10× {[}3.07, 5.80{]} ★} \\
\end{longtable}
}

Los tres enunciados sobre retrieval-como-palanca-dominante (§5.9.4 en versión
EN) se sostienen sobre los cuatro SoCs en dos OSes al mismo nivel de
significancia. La reproducción también confirma que \textbf{las conclusiones de
§5.9 son claims sobre el storage engine, no sobre el LLM runtime ni sobre el OS
runtime}: los chips difieren \textasciitilde11× en throughput LLM bruto (p50 de
\texttt{large\ +\ RAG}: iPhone 12 Pro 30.84 s, T760 49.23 s, SD662 71.99 s,
Kirin 659 152.26 s) pero las ratios de F1 que sostienen la tesis de §5.9 son
estables. La fila del Kirin también confirma la tesis de storage-engine desde la
dirección opuesta: retrieval (scan escalar FLAT @ \textasciitilde2.5 ms p50 para
N = 2 000, dim = 384) e ingest (addBatchDirect 47 ms para N = 2 000) operan
sobre los presupuestos default del SDK incluso sobre un chip Cortex-A53 / 4 GB
del 2017, así que la surface del engine documentada en este paper es portable
sobre las tres generaciones de microarquitectura Android testeadas (A76, A73,
A53).

\textbf{Descomposición de latencia.} Partir el \texttt{total\_us\ p50} por query
en sus cuatro componentes del pipeline (\texttt{embed} = forward de BGE sobre la
pregunta, \texttt{search} = retrieval vector, \texttt{prefill} = LLM ingiere
prompt + passages retrieved, \texttt{decode} = LLM autoregresivo hasta
\texttt{max\_new\_tokens\ =\ 64}) afina la tesis de storage-engine a un
statement de presupuesto:

\textbf{Tabla 18 --- small + RAG descomposición de latencia (mediana por query,
ms). Prefill domina, retrieval es \textless{} 0.13 \% del total en cada chip
incluyendo Cortex-A53 del 2017 y Apple A14 del 2020; la ratio
\texttt{prefill\ ≈\ 6×\ decode} que el perfil compute-vs-memory de Qwen 2.5 fija
es invariante a través de microarquitecturas \emph{y} sistemas operativos.}

{\def\LTcaptype{none} % do not increment counter
\begin{longtable}[]{@{}llllll@{}}
\toprule\noalign{}
Chip & embed p50 & search p50 & prefill p50 & decode p50 & total p50 \\
\midrule\noalign{}
\endhead
\bottomrule\noalign{}
\endlastfoot
Apple A14 (iPhone 12 Pro) & 14.2 & 1.7 & 11 835 & 2 016 & 13 576 \\
Unisoc T760 & 21.5 & 0.6 & 15 027 & 2 924 & 17 618 \\
QCOM SD662 & 30.8 & 0.9 & 22 237 & 4 387 & 26 237 \\
HiSi Kirin 659 & 67.4 & 3.7 & 48 374 & 9 138 & 56 159 \\
\end{longtable}
}

Dos invariantes visibles sobre los tres chips:

\begin{enumerate}
\def\labelenumi{\arabic{enumi}.}
\item
  \textbf{\texttt{prefill\ ≈\ 6×\ decode}.} El prefill compute-bound (cada token
  lee el KV cache completo) y el decode memory-bound (cada token streamea una
  vez) sientan en la misma ratio de FLOPs por token que Qwen 2.5 especifica, sin
  importar la microarquitectura. T760 es 3.2× más rápido que Kirin en términos
  absolutos, pero el split intra-query es el mismo.
\item
  \textbf{\texttt{embed\ +\ search\ ≪\ 1\ \%\ de}total\_us``} en cada celda:

  {\def\LTcaptype{none} % do not increment counter
  \begin{longtable}[]{@{}llllll@{}}
  \toprule\noalign{}
  Chip & embed & search & prefill & decode & retrieval / total \\
  \midrule\noalign{}
  \endhead
  \bottomrule\noalign{}
  \endlastfoot
  Apple A14 (iPhone 12 Pro) & 0.10 \% & 0.012 \% & 87.2 \% & 14.9 \% &
  \textbf{0.12 \%} \\
  Unisoc T760 & 0.12 \% & 0.003 \% & 85.3 \% & 16.6 \% & \textbf{0.13 \%} \\
  QCOM SD662 & 0.12 \% & 0.004 \% & 84.8 \% & 16.7 \% & \textbf{0.12 \%} \\
  HiSi Kirin 659 & 0.12 \% & 0.007 \% & 86.1 \% & 16.3 \% & \textbf{0.13 \%} \\
  \end{longtable}
  }

  La tesis de storage-engine de §5.9 colapsa a un statement de presupuesto, no a
  un argumento de tuning: una regresión de un orden de magnitud en el path de
  retrieval (e.g.~bajar el recall de HNSW a 80 \% o correr FLAT brute-force a 5×
  los ms que cuesta en Kirin) todavía dejaría el retrieval bajo 1.5 \% del
  total, bien dentro del noise floor del prefill. La spread 3× wall-clock entre
  chips es 100 \% LLM-runtime; el engine documentado en este paper ni ayuda ni
  perjudica ese envelope.
\end{enumerate}

\textbf{Techo de recall del retrieval.} Una pregunta natural a la descomposición
de latencia es si el pipeline de retrieval es el bottleneck del F1 de §5.9.
Computar \texttt{recall@k} de BGE-small sobre las mismas 200 queries NQ contra
el mismo corpus de 2 000 passages --- scoreando ``el gold passage de
short-answer está en el top-k'' contra el campo \texttt{gold} de cada query ---
da el F1 upper bound que el storage engine \emph{podría} entregar al LLM:

\textbf{Tabla 19 --- recall@k de retrieval de BGE-small sobre el slice de 200
queries NQ. Cosine sobre embeddings unit-normalisados de \texttt{dim\ =\ 384}.
Un solo número por row porque los embeddings son determinísticos y el chip
elegido no cambia el recall --- mismo modelo, mismo corpus, mismas queries.}

{\def\LTcaptype{none} % do not increment counter
\begin{longtable}[]{@{}ll@{}}
\toprule\noalign{}
k & recall@k \\
\midrule\noalign{}
\endhead
\bottomrule\noalign{}
\endlastfoot
1 & 0.905 \\
3 & 0.970 \\
5 & 0.980 \emph{(paper config)} \\
10 & 0.990 \\
\end{longtable}
}

El pipeline de retrieval entonces entrega el gold passage en el top-5 en 98 \%
de las queries --- pero \texttt{F1\_short} aterriza en 0.487 para
\texttt{large\ +\ RAG} sobre el chip más fuerte. El gap de 49 puntos entre
``retrieval perfecto'' y ``el modelo escribe la respuesta correcta'' es
\textbf{la quality de extraction del LLM, no el recall del storage engine}. El
small Qwen 2.5 0.5B compone esto --- \texttt{small\ +\ RAG} aterriza en 0.236,
la mitad del número del large-model, sobre el mismo contexto de 0.98 recall. Un
modelo class-7B cerraría la mayor parte del gap remanente sin cambiar una línea
del código del engine en este paper; la surface del storage engine satura su
upper bound sobre este corpus y ships cero budget bloqueante sobre F1.

\begin{quote}
\textbf{Sidebar de ingeniería --- fixes de portabilidad shipped durante este
sweep cross-platform.} Cinco issues SDK-level surgieron al llevar el bench más
allá del target original Moto G35 5G. Cada uno está documentado en
\texttt{research/results/cross\_platform\_e2e/} y el fix correspondiente se
shipped en la rama \texttt{kirin-4gb-sdk-opts}:

\begin{enumerate}
\def\labelenumi{\arabic{enumi}.}
\tightlist
\item
  \textbf{\texttt{flash\_attn} requiere \texttt{asimdhp}.} El path
  flash-attention de llama.cpp emula conversión fp16↔fp32 sobre cores sin
  half-precision nativa (Cortex-A53 / A73 baseline v8.0), lo que hace el path
  más lento y usa \emph{más} working memory que el kernel estándar. Cambiar el
  default de \texttt{flashAttention} a \texttt{CpuFeatures.hasFp16()} en lugar
  de \texttt{true} incondicional mantiene el path lento fuera de chips que no
  pueden usarlo.
\item
  \textbf{\texttt{n\_batch} del embedder capeado bajo \texttt{n\_ctx} divide
  passages largos en prefills multi-batch, lo que deadlockea en cores v8.0
  dentro del fallback fp16 de ggml.} El slice NQ de §5.9 tiene al menos un
  passage de \textasciitilde450 tokens (\texttt{passages{[}2{]}}); el default
  SDK \texttt{n\_batch\ =\ min(n\_ctx,\ 256)} para el embedder dividía ese
  prefill en 256 + 194 sub-batches y se congelaba en Kirin 659 / SD662 v8.0.
  Subir el default a \texttt{n\_batch\ =\ n\_ctx} (= 512) evita el path split en
  cada passage del slice.
\item
  \textbf{El paralelismo de \texttt{addBatchDirect} bulk no auto-throttle.} El
  pool default de 8-way \texttt{std::thread} deadlockea bajo CPU cgroup
  throttling agresivo (EMUI iAware), donde el kernel no puede mantener calientes
  los 8 workers y los threads spinean sobre el mutex per-element de hnswlib.
  Exponer \texttt{VectorIndex.setAddBatchThreads(n)} y un env var
  \texttt{DAZZLE\_HNSW\_BATCH\_THREADS} matching permite que un device tight-RAM
  pin el build pool a un solo worker sin cambiar parámetros del paper.
\item
  \textbf{La importance del foreground-service notification debe ser ≥ HIGH en
  EMUI 9.} El channel previo \texttt{IMPORTANCE\_LOW} causaba que el bench
  process se demote a \texttt{WORKINGSET\_BACKGROUND} (subCmd 352 de iAware)
  \textasciitilde10 s después del foreground grant sin importar wakelock o
  whitelist de batería; subir la importance del channel y agregar un
  re-\texttt{notify} heartbeat de 4 s es el workaround diagnosticado.
\item
  \textbf{Runs lanzados por activity en EMUI 9 requieren
  \texttt{am\ start\ -a\ MAIN\ -n} (la forma con \texttt{-c\ LAUNCHER} re-routea
  via HwLauncher y descarta extras del intent).} Runs lanzados por
  instrumentation pasan por encima del iAware throttling completamente; el
  harness ahora ships un JUnit entry point \texttt{RagE2EBenchTest} invocable
  via \texttt{am\ instrument}, con la misma surface de extras del intent que la
  activity lee.
\item
  \textbf{HNSW \texttt{addPoint(label\ =\ 0)} deadlockea sobre Cortex-A53 +
  libstdc++ Bionic.} La primera llamada a
  \texttt{hnswlib::HierarchicalNSW::addPoint} bloquea indefinidamente en Kirin
  659 / EMUI 9 / kernel 4.9 incluso con un \texttt{addBatchDirect} single-thread
  y un load mínimo de vectores random (4 vecs, dim = 4), confirmando que el bug
  está en la lib, no en el bench. \texttt{Algorithm.FLAT} (BruteforceSearch) es
  exacto y termina en 47 ms sobre N = 2 000, dim = 384 --- la fila del Kirin de
  la Tabla 17 reporta bajo FLAT. La causa raíz del deadlock está documentada en
  ocho passes de investigación; el SDK va a hacer fallback automático a FLAT
  sobre chips que matchean el fingerprint en v3.
\item
  \textbf{El kill-score de iAware en EMUI 9 acumula a través del embed loop y
  dispara sobre el siguiente mmap grande.} Incluso con el workaround FLAT
  shipped, el bench process es killed en silencio en \texttt{DazzleLlm.open}
  después de un embed loop de 25 min exitoso --- sin importar \texttt{useMmap},
  \texttt{useMlock}, \texttt{n\_threads} o pre-warm. Un probe standalone de
  instrumentation en un proceso fresh abre el mismo Qwen 0.5B en 1.5 s y exit
  limpio, confirmando que el kill es el score per-process de iAware, no un
  límite OS-wide del hardware. El harness ships un driver multi-process de
  cuatro fases (\texttt{RagE2EBenchPhases}) que parte el bench en invocaciones
  separadas de \texttt{am\ instrument}: \texttt{phase=embed} escribe embeddings
  a un cache binary, \texttt{phase=small} corre las variantes Qwen 0.5B en un
  proceso fresh (variant A + variant C back-to-back), y \texttt{phase=large}
  corre las variantes Qwen 1.5B en otro.
\item
  \textbf{El path FLAT del SDK mal-conectado a través de
  \texttt{BruteforceSearch} de hnswlib, que segfaultea sobre NEON Cortex-A53.}
  Con (7) shippeado el bench todavía obtenía 0.025 en \texttt{small\ +\ RAG} (vs
  0.235 en los chips HNSW) porque \texttt{RagE2EBenchPhases.ensureServerRunning}
  arrancaba Valkey sin \texttt{DazzleModule.VectorSearch}, así que
  \texttt{FT.CREATE} no tenía handler y el schema per-index nunca aterrizaba en
  \texttt{g\_indexes} --- \texttt{addBatchDirect} / \texttt{searchDirect}
  JNI-direct retornaban en silencio con 0 hits. Con el módulo cargado, la
  siguiente capa del bug emergió: el SIMD distance kernel de
  \texttt{BruteforceSearch::searchKnn} de hnswlib segfaultea sobre loads NEON
  16-byte misaligned sobre el stride packed \texttt{{[}vec,\ label{]}} que
  hnswlib usa, en Cortex-A53 + libstdc++ Bionic. El fix mirrorea cada vector del
  path FLAT en \texttt{schema-\textgreater{}fp32\_store} y reemplaza
  \texttt{searchKnn} con un scan brute-force escalar portable;
  \textasciitilde2.5 ms por query sobre N = 2 000, dim = 384 en Cortex-A53. El
  path HNSW no cambia.
\item
  \textbf{Split-prefill del LLM deadlockea en Cortex-A53 v8.0 cuando el largo
  del prompt excede \texttt{n\_batch}.} Con (7) y (8) shipped, el bench pasaba
  las fases de embed y addBatch pero el bench process era killed dentro de
  segundos de la primera query \texttt{+RAG} en \texttt{DazzleLlm.generate}. Las
  variantes \texttt{no-RAG} completaban --- distinto code path, prompt de
  \textasciitilde30 tokens. El prompt \texttt{+RAG} es de \textasciitilde570
  tokens y \texttt{cfg.llmNBatch\ =\ 512}, así que llama.cpp partía el prefill
  en una primera call de 512 tokens + una continuación de 58 tokens, golpeando
  el mismo fallback fp16 ggml v8.0 que el item 2 documentó para el embedder.
  Subir \texttt{llmNBatch} a \texttt{llmNCtx} (= 2048) hace prefill de cualquier
  prompt hasta el context size en una sola call; las queries \texttt{+RAG}
  entonces completan con \texttt{prompt\_tokens.avg\ =\ 570} matcheando T760 /
  SD662.
\item
  \textbf{iAware re-dispara después de \textasciitilde30 queries de
  \texttt{Qwen-1.5B\ +\ RAG} incluso con los items 7--9 aplicados.} La variante
  D (large + RAG) consistentemente muere entre query 5 y query 30 a través de
  reboots y reinstalaciones. El harness del bench ahora ships un runner chunked
  de variantes (\texttt{runRagE2EVariantChunk} con extras \texttt{q\_offset} /
  \texttt{q\_limit}) para que cada chunk corra en un proceso
  \texttt{am\ instrument} fresh y los archivos JSON per-chunk sean concatenados
  por la fase de merge. La fila del Kirin de la Tabla 17 fue producida por 20
  chunks de 10 queries para variant D más 4 chunks de 50 queries para variant B;
  wall clock total para Phase 2b ≈ 8 hr 50 min.
\item
  \textbf{Watchdog de launch iOS (\texttt{0x8BADF00D}) en el primer run de
  RagE2EBench.} El bench monolítico \texttt{RagE2EBench.run()} invocado inline
  desde \texttt{DazzleExperimentApp.init()} agota el watchdog de launch de 20
  segundos de iOS antes de que SwiftUI tenga chance de renderizar, así que
  FrontBoard mata el proceso con
  \texttt{0x8BADF00D\ ProcessVisibility:Foreground}. Fix: dispatchear el bench
  en \texttt{DispatchQueue.global(qos:\ .userInitiated).async} desde
  \texttt{init()} para que el main runloop de SwiftUI agende la view placeholder
  dentro de la ventana del watchdog. El bench termina \textasciitilde25--35 min
  después en background queue y se hace \texttt{exit(0)} solo. Mismo fix lands
  como utility public en el entry point Swift de \texttt{RagE2EBench.run()} así
  cualquier app iOS que use el entry point obtiene el dispatch automático.
\item
  \textbf{Default \texttt{n\_batch\ \textless{}\ n\_ctx} del context LLM en iOS
  gatilla split prefill en prompts largos.} Con (11) shippeado, el bench iPhone
  pasaba embed + addBatchDirect pero el proceso era killed dentro de 1 s de la
  primera query \texttt{+RAG} en \texttt{llama\_decode} --- mismo pattern que el
  pass 15 del Kirin documentó para Android v8.0, esta vez en el build de
  llama.cpp iOS. \texttt{dazzle\_llama\_new\_context} shippeaba con
  \texttt{n\_batch\ =\ 512} hard-coded; el prompt de §5.9 es \textasciitilde570
  tokens así que \texttt{llama\_decode} parte el prefill en 512 + 58 sub-batches
  y el continuation aborta. Fix: setear \texttt{n\_batch\ =\ n\_ctx} en
  \texttt{dazzle\_llama\_new\_context} para que cualquier prompt hasta context
  size hace prefill en una sola llamada \texttt{llama\_decode}. El cambio aplica
  en cada plataforma; cuesta \textasciitilde30 MB extra de compute buffer a
  \texttt{n\_ctx\ =\ 2048}, despreciable vs los pesos del modelo.
\item
  \textbf{Método convenience Swift \texttt{DazzleServer.vectorIndex(...)} sin
  \texttt{initialCapacity}.} Con (11) y (12) shippeados, el bench pasaba la fase
  LLM pero \texttt{addBatchDirect} abortaba en el elemento 1024 con
  \texttt{BruteforceSearch::addPoint\ runtime\_error("exceeds\ limit")} de
  hnswlib porque el index FLAT se creaba con el default SDK
  \texttt{INITIAL\_CAP\ =\ 1024}. El bench Kotlin Android pasaba 2 000 vía el
  ctor low-level \texttt{vectorIndex(...)} que acepta \texttt{initialCapacity};
  el método convenience Swift en \texttt{DazzleServer.shared.vectorIndex(...)}
  no lo exponía. Fix: extender la firma \texttt{DazzleServer.vectorIndex} iOS
  con \texttt{initialCapacity:\ Int\ =\ 0} (más \texttt{m},
  \texttt{efConstruction}) para que los callers puedan pre-sizear como en el
  lado Android.
\end{enumerate}

Los items 1--5 son fixes universales de portabilidad que landed sobre los cuatro
chips. Los items 6--10 son específicos del path Kirin 659 y los items 11--13
específicos del port iOS; ninguno cambia la comparación numérica de las celdas
de la Tabla 17. Juntos cierran el gap entre ``el bench corrió sobre el chip
donde lo desarrollamos'' y ``el bench corre sobre un spread de SoCs Android
mid-range y un SoC Apple sobre dos sistemas operativos'' --- una precondición
para cualquier claim cross-platform / cross-OS sobre una superficie de engine
embebida.
\end{quote}

\section{Discusión}\label{discusiuxf3n}

\subsection{Agregados materializados como primitiva, no como
optimización}\label{agregados-materializados-como-primitiva-no-como-optimizaciuxf3n}

El retrieval de baja latencia en este paper monta sobre el patrón de agregados
materializados, no sobre alguna marca específica de motor de base de datos. Ya
sea implementado como un hash mantenido por Lua en Dazzle, una tabla mantenida
por triggers en SQLite, o un step explícito de update en una estructura
in-memory, el patrón desplaza el trabajo de agregación al write time y hace
bounded el retrieval en read time.

La contribución de Dazzle no es inventar este patrón. La contribución es
entregarlo como una superficie primitiva first-class
(\texttt{EVALSHA}-maintained hash + \texttt{HMGET}) con acceso SDK tipado, así
que el autor de la aplicación no necesita construir lógica de migración,
orquestación de triggers, o código de mantenimiento manual para cada app.

Esta distinción importa para la tesis del paper: en el envelope de latencia
donde todos los backends son aceptables, el factor decisivo es menos ``¿puede
existir el patrón?'' y más ``cuánta ceremonia de ingeniería se requiere para
hacerlo correcto y mantenible en producción''.

\subsection{El bloque de contexto bounded}\label{el-bloque-de-contexto-bounded}

La propiedad de \textasciitilde58--62 tokens independiente de N viene de una
decisión arquitectónica, no de coincidencia. El bloque de contexto codifica seis
escalares agregados (min/max/avg/count/\ldots) más el sumario probabilístico, no
una lista de N lecturas crudas. El tamaño del prompt por tanto se desacopla del
tamaño del dataset: un agente corriendo durante semanas en un dispositivo no
debería consumir progresivamente más de la ventana de contexto del modelo
simplemente porque más datos se han acumulado.

\subsection{Limitaciones}\label{limitaciones}

\begin{itemize}
\item
  \textbf{Superficie primitiva contada como API nativa, no como equivalencia
  expresiva.} La Tabla 1 (§2.2) reporta disponibilidad de primitivas como un
  check categórico (\texttt{N} API nativa / \texttt{E} extensión oficial /
  \texttt{A} re-implementación a nivel de aplicación / \texttt{T} extensión
  third-party) sobre cómo el backend llega a cada construcción, con el conteo
  headline medido solo sobre \texttt{N}. Las primitivas marcadas \texttt{E},
  \texttt{A} o \texttt{T} para backends no-Dazzle son \emph{implementables}
  atómicamente a través de las facilidades general-purpose del motor subyacente
  --- por ejemplo, range queries de sorted-set vía \texttt{ORDER\ BY} +
  \texttt{WHERE\ BETWEEN} en SQLite, increments atómicos de float vía
  \texttt{UPDATE\ …\ SET\ v\ =\ v\ +\ ?} dentro de una transacción, indexación
  geográfica vía la extensión \texttt{R*Tree} de SQLite (extensión oficial
  enviada con la SQLite amalgamation desde 2008), bounded \texttt{MAXLEN}
  streams vía triggers, y TTL vía un job \texttt{DELETE} programado. La señal de
  la Tabla 1 es por tanto ``cuánto código pegamento escribe el autor del agente
  para alcanzar esta primitiva'' (una cuestión ergonómica de superficie
  primitiva que ata de vuelta a la Tabla 9), no ``esta funcionalidad es
  alcanzable en este motor'', que en casi cada caso, sí.
\item
  \textbf{Post-mortem del harness.} Un bug en la truth-source
  \texttt{SqliteBruteforceVector} causó que data obsoleta se acumulara a través
  de configuraciones del benchmark recall-floor en v1; fue arreglado en el
  commit \texttt{1e3d5f5} y cada celda de las Tablas 4, 11 y 12 de esta revisión
  viene del run post-fix. La timeline completa (introducción, detección, fix,
  asserts que atrapan una regresión futura, y una auditoría de cada otro
  benchmark on-device en el repo por el mismo bug pattern) está documentada en
  el Apéndice B (\texttt{research/paper/appendix\_b\_harness\_postmortem.md}).
\item
  Evaluación end-to-end LLM en un solo dispositivo físico (Moto G35 5G) +
  validación cross-platform en iPhone 12 Pro. Los resultados en otros SoCs y
  versiones OS pueden diferir; los reportamos en tres familias adicionales de
  SoC Android (Unisoc, MediaTek, Qualcomm) en companion benchmark reports.
\item
  Carga sintética. Streams reales de IoT tienen patrones bursty, fallos de
  sensor, y concurrencia adicional que no modelamos.
\item
  La evaluación de accuracy del modelo usa tres estadísticas verificables; una
  evaluación más rigurosa usaría un set de preguntas factuales más grande y
  múltiples runs por condición.
\item
  Paridad de transporte iOS/Android incompleta: el SPSC ring buffer y el path
  \texttt{io\_uring} son Linux-only; iOS se queda sobre el Phase-0 pipe path,
  suficiente para la carga evaluada (§5.5).
\item
  \textbf{Carga LLM end-to-end de un solo dispositivo.} El bench end-to-end RAG
  de §5.9 se reporta en un solo dispositivo físico (Moto G35 5G).
\item
  \textbf{Los peers comerciales evolucionan.} SQLiteAI v0.9.95 no expone HNSW al
  cierre de este bench; una versión futura del producto que incluya HNSW
  desplazaría la comparación de §5.8 de ``HNSW vs brute-force scan'' a ``HNSW vs
  HNSW'', reduciendo el gap a constant factors en lugar de asimetría
  algorítmica.
\item
  \textbf{Skew de footprint allocator-platform.} El tamaño on-disk e in-memory
  de un backend embebido es sensible al allocator host y a la configuración
  default. Observamos un delta cross-platform de 23× sobre Dazzle (153 KB en iOS
  vs 6.7 KB en Android, mismo Valkey-fork in-process, misma carga de 200
  lecturas) atribuible al \texttt{libsystem\ malloc} de Apple versus el
  comportamiento de la arena \texttt{jemalloc} de Android, y un delta de 42×
  sobre RocksDB en la dirección opuesta (4.41 MB Android vs 104 KB iOS)
  impulsado por los defaults de pre-allocation WAL + SST de Android.
\item
  \textbf{Comparación de configuración default en §5.2.} La comparación
  storage-latency usa configuraciones de backend default. Un setup SQLite con
  vistas materializadas mantenidas manualmente vía triggers cerraría la mayor
  parte del gap medido. La contribución de Dazzle es shipear el patrón
  nativamente, no claimar superioridad algorítmica sobre SQLite.
\item
  \textbf{El benchmark vectorial abarca dos clases algorítmicas.} En §5.8,
  Dazzle y ObjectBox corren HNSW (O(log N)), mientras que SQLiteAI v0.9.95
  expone brute-force quantized scan (O(N)). Una release futura de SQLiteAI con
  HNSW user-accessible desplazaría la comparación a constant factors.
\item
  \textbf{El benchmark de aplicación RAG aísla disponibilidad de retrieval, no
  identidad de motor vectorial.} La sección §5.9 compara with-retrieval vs
  without-retrieval y no aísla el motor vectorial de Dazzle contra ObjectBox o
  SQLiteAI en la misma pipeline end-to-end.
\item
  \textbf{Sesgo de autor sobre la medición LOC.} El autor de este paper es
  también el autor de las seis implementaciones de backend contra las que Dazzle
  se mide en la Tabla 9. Los managers Dazzle-side fueron los primeros escritos y
  los más iterados; los managers SQLite/LMDB/RocksDB/ObjectBox son ports
  directos de la misma superficie de capa de estado del agente pero no fueron
  optimizados subsecuentemente por line count.
\item
  \textbf{Conflicto de interés y compromiso público.} El autor de este paper es
  también el autor de \texttt{dazzle-sdk} (Apache-2.0; Maven Central / pub.dev /
  npm / Swift Package Manager) y de cada wrapper de comparison-engine
  benchmarkado aquí. Para hacer la comparación falsificable en lugar de
  auto-servida, tres compromisos concretos sostienen desde la fecha de esta
  revisión en adelante:

  \begin{enumerate}
  \def\labelenumi{\arabic{enumi}.}
  \tightlist
  \item
    \textbf{Los seis wrappers de backend son artefactos first-class en el
    repositorio.} Todo su source, build configuration, y test entry points viven
    bajo \texttt{experiment/backends/} y tienen la misma licencia Apache-2.0 que
    Dazzle. Cualquiera puede forkear, re-tunear, y re-correr.
  \item
    \textbf{External pull requests que reduzcan el LOC count de cualquier
    wrapper no-Dazzle, o que mejoren cualquier número medido del backend
    no-Dazzle en las Tablas 3 / 5 / 11, serán revisados y mergeados sobre mérito
    técnico.}
  \item
    \textbf{Cada optimización aceptada se registra en
    \texttt{experiment/backends/CHANGELOG.md} antes del próximo bench freeze.}
  \end{enumerate}

  La metodología del benchmark, los outputs JSON crudos, y los scripts de
  análisis están abiertos en el repositorio (paths listados en §8.1
  Reproducibility) así que cualquier claim en el paper puede re-correrse contra
  un build independiente de los motores de comparación.
\end{itemize}

\subsubsection{Validación cross-platform del benchmark
vectorial}\label{validaciuxf3n-cross-platform-del-benchmark-vectorial}

Las Tablas 11 y 12 se reportan en un solo dispositivo físico (Moto G35 5G,
Unisoc T760, ARMv8.2-A). Para descartar artefactos single-device, re-corrimos el
mismo harness sobre tres familias adicionales de SoC Android y re-bootstrapeamos
la celda headline N = 20 000 con paired-query resampling
(\texttt{B\ =\ 10\ 000}, seed = 42). Las tablas per-engine completas están en
\texttt{research/paper/vecbench\_cross\_platform\_ci.md}; los headline ratios se
reproducen aquí:

{\def\LTcaptype{none} % do not increment counter
\begin{longtable}[]{@{}
  >{\raggedright\arraybackslash}p{(\linewidth - 8\tabcolsep) * \real{0.2500}}
  >{\raggedright\arraybackslash}p{(\linewidth - 8\tabcolsep) * \real{0.3171}}
  >{\raggedleft\arraybackslash}p{(\linewidth - 8\tabcolsep) * \real{0.1585}}
  >{\raggedleft\arraybackslash}p{(\linewidth - 8\tabcolsep) * \real{0.1890}}
  >{\raggedright\arraybackslash}p{(\linewidth - 8\tabcolsep) * \real{0.0854}}@{}}
\toprule\noalign{}
\begin{minipage}[b]{\linewidth}\raggedright
Device
\end{minipage} & \begin{minipage}[b]{\linewidth}\raggedright
SoC (ISA, big core)
\end{minipage} & \begin{minipage}[b]{\linewidth}\raggedleft
Round 1 p50 (dispatched)
\end{minipage} & \begin{minipage}[b]{\linewidth}\raggedleft
Round 2 p50 (baseline-forced)
\end{minipage} & \begin{minipage}[b]{\linewidth}\raggedright
Δ
\end{minipage} \\
\midrule\noalign{}
\endhead
\bottomrule\noalign{}
\endlastfoot
Moto G35 5G \emph{(reference)} & Unisoc T760 (ARMv8.2-A + fp16/dot, Cortex-A76)
& 269 {[}261, 273{]} µs & 269 {[}264, 273{]} µs & \textbf{0 µs} \\
Huawei Y9a (FRL-L23) & MediaTek Helio G80 (ARMv8.2-A + fp16/dot, A75) & 671
{[}642, 701{]} µs & 588 {[}566, 612{]} µs & \textbf{-83 µs} ² \\
Moto G30 & Snapdragon 662 (ARMv8.0-A, Cortex-A73) & 445 {[}409, 478{]} µs & 519
{[}488, 548{]} µs & +74 µs ¹ \\
Huawei P20 Lite (ANE-LX3) & Kirin 659 (ARMv8.0-A, Linux 4.9, Cortex-A53) & 1054
{[}1035, 1092{]} µs & 1043 {[}1014, 1078{]} µs & -11 µs ¹ \\
\end{longtable}
}

¹ Los chips ARMv8.0 corren la misma \texttt{libdazzle.so} baseline en ambas
rondas (la variante v82 haría SIGILL); el Δ ronda-a-ronda es por tanto puro
ruido térmico/load, ranging ±10 \% en N = 20 000.

² En Helio G80 (Cortex-A75 {[}@armcortexa75{]} big core) el binario v82
dispatched es \textbf{más lento} que el baseline en ≈ 12 \% en p50 (los CI 95 \%
no se solapan: {[}642, 701{]} vs {[}566, 612{]} µs). Los kernels NEON del
hot-path son idénticos --- \texttt{simsimd} runtime-selecciona los mismos
símbolos \texttt{\_neon\_f16\ /\ \_neon\_dotprod} en ambos binarios --- así que
el gap está en el código C++ circundante.

\subsection{Asimetría algorítmica en el benchmark
vectorial}\label{asimetruxeda-algoruxedtmica-en-el-benchmark-vectorial}

La comparación en §5.8 abarca dos clases algorítmicas: HNSW (Dazzle, ObjectBox)
es approximate nearest neighbour search con expected \(O(\log N)\) search
complexity {[}@hnsw{]}; SQLiteAI v0.9.95 implementa linear scan acelerado por
SIMD \(O(N)\) sobre vectores cuantizados int8. La asimetría es por diseño del
producto SQLiteAI --- \texttt{vector\_quantize\_scan} es el único query path
optimizado que el producto expone; HNSW no es user-accessible en esta versión.

Dentro de la Tabla 11 la lectura de operating-point es la que se debe usar:
\textbf{Dazzle SQ8 aterriza en recall = 0.959 / p50 = 208 µs, ObjectBox 4.x en
recall = 0.994 / p50 = 853 µs, SQLiteAI precompute en recall = 0.987 / p50 = 1
407 µs}, todos en N = 20 000, dim = 384, recall floor \(\geq 0.95\).

Los dos motores HNSW (Dazzle SQ8 y ObjectBox) son \textbf{dos puntos sobre la
misma curva Pareto recall--latency}, no winner / loser: Dazzle SQ8 trades
precisión de cuantización int8 por p50 \textasciitilde4× menor; ObjectBox
mantiene precisión fp32 y paga el costo de latencia. \textbf{Ambas elecciones se
mantienen dentro del envelope LLM-noise-floor en este N}. SQLiteAI se ubica en
una clase algorítmica distinta (\(O(N)\) scan); su celda 1 407 µs todavía cabe
en el noise floor en N = 20 000 pero no en N \(\gg 20\,000\) --- el threshold de
tamaño de corpus que §5.8.4 documenta.

\section{Trabajo relacionado}\label{trabajo-relacionado}

\textbf{Inferencia on-device.} LiteRT-LM, llama.cpp, MLC-LLM, y ExecuTorch
{[}@litertlm; @llamacpp; @mlcllm; @executorch{]} apuntan a compresión de modelo
y eficiencia de runtime. Ninguno aborda la capa de ejecución stateful; cada
inference call se trata como un evento independiente.

\textbf{Memoria para agentes.} MemGPT {[}@memgpt{]} propone una jerarquía de
memoria de dos niveles (in-context + external) targeting server deployment.
Generative Agents {[}@generativeagents{]} usa un memory stream
retrieval-augmented para simulación NPC. A-MEM {[}@amem{]} introduce memoria
agent-managed con linking dinámico. Las tres asumen entornos cloud o server con
un proceso de base de datos separado, red, y RAM substancial. Dazzle mueve esa
capa de memoria externa \textbf{dentro del proceso de la app móvil}.

\textbf{Redis sobre hardware restringido.} Redis sobre Raspberry Pi alcanza
latencia sub-milisegundo pero corre como demonio separado sobre TCP. No
conocemos un Valkey-on-Android port publicado.

\textbf{Motores embebidos.} SQLite, LMDB, y RocksDB están diseñados desde cero
como librerías in-process. Dazzle toma el path opuesto: arranca de un motor
server diseñado alrededor de TCP loopback y un event loop single-threaded, y
reescribe su substrato de I/O (no sus estructuras de datos) para hacer ejecución
in-process el default. El resultado se sienta entre ambos campos. La biblioteca
se consume exactamente como un motor embebido, mientras que el modelo de datos
es el que un full server-grade database expone.

\textbf{Bases de datos vectoriales móviles --- tres productos distintos.} El
ecosistema mobile SQLite-vector incluye tres productos comúnmente confundidos
que vale la pena distinguir aquí. (1) \textbf{SQLite} plano no tiene soporte
nativo para vectores y se benchmarkea solo sobre su capa de storage en §5.2. (2)
El \textbf{sqlite-vec} open-source de Alex Garcia
(\texttt{github.com/asg017/sqlite-vec}, Apache 2.0) es una extensión SQLite
agregando vector search brute-force; esta revisión lo benchmarkea en el sweep
SQLite-family (§5.8.4). (3) \textbf{SQLiteAI sqlite-vector}
{[}@sqliteaivector{]} (\texttt{github.com/sqliteai/sqlite-vector}, Elastic
Licence 2.0) es el producto comercial de SQLite Cloud, Inc.~(la entidad
corporativa detrás del ecosistema SQLite AI); ship un path de query linear-scan
cuantizado acelerado por SIMD (\texttt{vector\_quantize\_scan}) y
deliberadamente evita HNSW. \textbf{ObjectBox 4.x} {[}@objectbox4{]} es el otro
peer comercial en nuestro benchmark y ship HNSW nativo respaldado por vectores
float32.

\textbf{Librerías ANN vectoriales server-grade.} \textbf{FAISS} {[}@faiss{]}
(Facebook AI), \textbf{DiskANN} {[}@diskann{]} (Microsoft Research), y
\textbf{Milvus} {[}@milvus{]} (Zilliz) fijan el baseline open-source para
approximate nearest-neighbour search a escala billion-point sobre hardware
server. La literatura de cuantización vectorial debajo de esos motores ---
Product Quantisation {[}@jegoupq{]} en particular --- es el precedente
algorítmico para el path SQ8 que Dazzle usa.

Los kernels NEON de Dazzle para fp32 dot product, fp32 L2-squared, fp16 dot
product, e i8 cosine vienen de \texttt{simsimd} {[}@simsimd{]} (Vardanian,
Apache 2.0 license) y usan el runtime dispatcher de esa librería
(\texttt{simsimd\_capabilities()} sobre \texttt{/proc/cpuinfo}) para seleccionar
la mejor variante disponible en cada chip.

\textbf{Object stores móviles más allá de SQLite.} \textbf{Realm} {[}@realm{]}
(ahora MongoDB) y \textbf{WCDB} {[}@wcdb{]} (Tencent / WeChat) son los dos peers
production-deployed de ObjectBox en el espacio móvil object-store. Ambos
targetean Android + iOS, ambos ship APIs typed más altas que SQLite raw, y ambos
son ampliamente usados en apps comerciales. Ninguno ofrece vector search o
HyperLogLog como primitiva (Realm ship full-text search; WCDB es
SQLite-on-the-wire underneath).

\textbf{Evaluación de RAG.} \textbf{RAGTruth} {[}@ragtruth{]} es el corpus
canónico de hallucination para evaluación retrieval-augmented LLM al cierre de
2024. La sección §5.9 usa NQ-open en su lugar porque la unidad de medición aquí
es \emph{recall de canonical short-answer strings} (si el LLM emite el gold span
cuando retrieval está presente). RAGTruth mide hallucination \emph{dado}
contexto retrieved, que es una señal estrictamente downstream que requiere un
hallucination judge. Layerar evaluación estilo RAGTruth sobre el setup de §5.9
es trabajo futuro; pertenecería al mismo eje que el reporting EM-vs-F1 que ya
incluimos.

\section{Conclusión}\label{conclusiuxf3n}

Dazzle propone una lente diferente de selección de backend para agentes
on-device. La superficie primitiva, la ergonomía de desarrollador, y el envelope
de latencia importan más que carreras aisladas en microsegundos. En la carga
medida cada backend evaluado opera dentro de un envelope aceptable de retrieval;
los diferenciadores operativos son qué primitivas de agent-state están
disponibles nativamente y cuánto esfuerzo de ingeniería cuestan de usar.

Tres resultados empíricos respaldan ese framing. La latencia de retrieval en el
path de estado estable de Dazzle es aproximadamente 0.00176 \% de un solo turno
de inferencia LLM on-device (50 µs contra \(\sim 2.84\) s para Gemma 4 E2B en el
Moto G35 5G), lo que ubica el costo de retrieval del backend dentro del piso de
ruido del bucle end-to-end. Superficie primitiva y esfuerzo de desarrollo
difieren materialmente: Dazzle expone diez primitivas nativas contra \(\leq 4\)
en las alternativas embebidas medidas, con implementaciones representativas de
agent-state aterrizando alrededor de 175 LOC (Android) / 186 LOC (iOS) en Dazzle
contra aproximadamente 200--290 LOC a través de las alternativas (Tabla 9, ambas
plataformas). A nivel de aplicación, una ablación 2×2 RAG eleva
\texttt{EM\_contains} 6.0× sobre Qwen 2.5 0.5B (0.105 → 0.630) y 6.7× sobre Qwen
2.5 1.5B (0.110 → 0.735); un modelo de 380 MB con retrieval supera a un modelo
3× más grande sin retrieval en cada métrica factual, end-to-end en un
dispositivo Android de \$150.

La pregunta práctica para el próximo deployment on-device por tanto se desplaza
de ``¿qué backend es más rápido?'' a ``¿qué set de primitivas y qué composición
de modelo encajan mejor en la aplicación, dado que el costo de retrieval es
efectivamente invisible al nivel de turno?''.

Dazzle se distribuye públicamente como \texttt{dazzle-sdk} 1.0.0-beta.4 sobre
Maven Central, pub.dev, npm, y Swift Package Manager. El código original de
Dazzle se publica bajo Apache-2.0; las porciones derivadas de Valkey retienen
BSD-3-Clause. Invitamos a la comunidad a validar estos resultados sobre hardware
móvil adicional y a extender la matriz del benchmark a través de cargas más
amplias.

\subsection{Artefacto y reproducibilidad}\label{artefacto-y-reproducibilidad}

Para hacer la evaluación auditable, este repositorio incluye los scripts,
dataset builders, y árboles de resultados crudos usados para producir las tablas
del paper. Los componentes están organizados así:

\begin{itemize}
\tightlist
\item
  \textbf{Bench runners}: \texttt{research/scripts/run\_experiment.sh},
  \texttt{run\_full\_benchmark.sh}, \texttt{run\_ios\_benchmark.sh},
  \texttt{run\_ablation\_sweep.sh},
  \texttt{run\_storage\_microbench\_per\_backend.sh},
  \texttt{run\_vector\_sqlite\_family.sh}.
\item
  \textbf{Generación de dataset y RAG slice tooling}:
  \texttt{research/scripts/generate\_dataset.py}, \texttt{nq\_slice.py},
  \texttt{recompute\_rag\_metrics.py}.
\item
  \textbf{Análisis de resultados y generación de tablas}:
  \texttt{analyze\_results.py}, \texttt{analyze\_storage\_microbench.py},
  \texttt{make\_vector\_bench\_table.py},
  \texttt{analyze\_vector\_sqlite\_family.py}.
\item
  \textbf{Árboles de resultados crudos y procesados}:
  \texttt{research/benchmarks/results/}, \texttt{research/results/}.
\end{itemize}

Cada métrica headline en §5 fue computada desde estos artefactos locales. Una
guía compacta paso-a-paso de reproducción se provee en
\texttt{research/paper/REPRODUCIBILITY.md} y puede ejecutarse end-to-end sobre
el mismo workspace layout.

\subsection{Política de versionado y compromiso público de
actualización}\label{poluxedtica-de-versionado-y-compromiso-puxfablico-de-actualizaciuxf3n}

Este paper documenta un moving target: \texttt{dazzle-sdk} es un artefacto vivo
bajo desarrollo activo a través de cuatro registries de paquetes, los backends
de comparación evolucionan, y varios hallazgos en esta revisión son
explícitamente null results que invitan a investigación adicional. Para mantener
al paper como una referencia útil en lugar de un snapshot congelado, el autor se
compromete a la siguiente política de versionado en arXiv:

\begin{enumerate}
\def\labelenumi{\arabic{enumi}.}
\tightlist
\item
  \textbf{Cada cambio medible en cualquier celda headline de las Tablas 3, 5,
  11, 12, o 15 produce una nueva revisión arXiv.}
\item
  \textbf{Cada nuevo dispositivo físico agregado a la tabla cross-platform
  (§6.3) produce una nueva revisión arXiv.}
\item
  \textbf{Cada PR externo aceptado que mejore los números medidos de un backend
  no-Dazzle produce una nueva revisión arXiv} --- con el PR contribuyente citado
  y los números viejos preservados con strikethrough.
\item
  \textbf{El abstract arXiv lleva una fecha de ``Last updated'' y un changelog
  de una línea del cambio material más reciente}.
\item
  \textbf{Todas las mediciones JSON crudas} viven en
  \texttt{research/benchmarks/results/} bajo paths timestamped, nunca
  sobrescritas a través de revisiones. Las re-derivaciones bootstrap-CI son
  deterministas (\texttt{B\ =\ 10\ 000}, \texttt{seed\ =\ 42}).
\item
  \textbf{El branch main del repositorio es el source of truth.} Cuando una
  discrepancia surja entre una revisión arXiv y \texttt{main}, \texttt{main}
  gana; la discrepancia gatilla una revisión por (1).
\end{enumerate}

El autor es alcanzable en \texttt{ivan.aliaga@urp.edu.pe} para preguntas de
revisor, issues de reproducción, y pull requests contra el código del benchmark.
Issues abiertos en el repositorio GitHub se acknowledgean dentro de siete días;
las revisiones arXiv numeradas se procesan dentro de catorce días del trigger
event.

\section{Referencias}\label{referencias}

Ver \texttt{research/paper/arxiv-build/refs.bib} para las entradas
bibliográficas completas. Esta versión en español preserva las citas verbatim
del paper en inglés (formato \texttt{{[}@key{]}} por entrada bibtex). 49
referencias numeradas.

\section{Apéndice A --- Capa de adaptadores LLM
(referencia)}\label{apuxe9ndice-a-capa-de-adaptadores-llm-referencia}

El SDK Dazzle expone una sola interfaz \texttt{LLMClient} (completion,
streaming, deltas de tool-call, lifecycle) y ship cinco adaptadores concretos
detrás de ella:

\begin{itemize}
\tightlist
\item
  \textbf{\texttt{LlamaCppClient}} --- adaptador local llama.cpp, GGUF q4\_k\_m
  / q5\_k\_m / q8\_0, n\_ctx configurable.
\item
  \textbf{\texttt{LiteRtLmClient}} --- adaptador LiteRT-LM (Google AI Edge),
  artefactos \texttt{.litertlm} cargados via runtime LiteRT.
\item
  \textbf{\texttt{FoundationModelsClient}} --- adaptador Apple Intelligence (iOS
  26+), enruta a Apple's on-device foundation model bridge.
\item
  \textbf{\texttt{OpenAICompatibleClient}} --- adaptador HTTP-compatible para
  llamar a APIs OpenAI-style (tanto endpoints OpenAI como servidores
  compatibles).
\item
  \textbf{\texttt{AnthropicClient}} --- adaptador para Anthropic Messages API
  {[}@anthropicapi{]}.
\end{itemize}

Las cinco implementaciones residen en
\texttt{sdk/android/src/main/java/dev/dazzle/sdk/edge/} (Kotlin) y
\texttt{sdk/ios/Sources-LiteRTLM/}, \texttt{sdk/ios/Sources-FoundationModels/}
etc. (Swift), comparten la misma interfaz tipada de
\texttt{completion(prompt,\ params)} /
\texttt{streamCompletion(prompt,\ params,\ onDelta)}, y son intercambiables
desde el agent-loop perspective.
